\documentclass[12pt,a4paper]{article}
\usepackage[utf8]{inputenc}
\usepackage[T1]{fontenc}
\usepackage{amsmath,amssymb,amsfonts}
\usepackage{mathtools}
\usepackage{booktabs}
\usepackage{array}
\usepackage{geometry}
\usepackage{hyperref}
\usepackage{cite}
\geometry{margin=2.5cm}

\newcommand{\MeV}{\,\text{MeV}}
\newcommand{\GeV}{\,\text{GeV}}
\newcommand{\fpi}{f_\pi}
\newcommand{\qqbar}{\langle\bar{q}q\rangle}
\newcommand{\zch}{\mathfrak{z}}

\title{A Tale of Two Tuples: Seed Triples, Signed Charges, \\
and SUSY Breaking in the sBootstrap}

\author{A.~Rivero\thanks{Manuscript prepared with AI assistance (Claude/Anthropic). 
Continues the program of \cite{Rivero2005,Rivero2012}.}}
\date{March 2026}

\begin{document}
\maketitle

\begin{abstract}
We reformulate the sBootstrap in terms of signed mass charges 
$\zch = z_0 + z_i$ with $m = \zch^2$, and show that the Koide ratio 
$Q = 2/3$ is the algebraic consequence of two conditions: 
$\sum z_i = 0$ and $\langle z_i^2\rangle = z_0^2$ (perturbation 
energies average to the base-scale energy). Allowing negative charges, 
we find that pseudoscalar mesons satisfy a Koide triple 
$(-\pi, D_s, B)$ to 0.1\%, and identify \emph{seed triples} with 
one massless member---$(0, s, c)$ for quarks, $(0, \pi, D_s)$ for 
mesons---as the pre-breaking configurations from which the full 
spectrum emerges. These seed triples share the same analytical 
structure (charge ratio $2-\sqrt{3}$) and, strikingly, the meson 
seed $(0,\pi,D_s)$ has Koide scale $v_0^2 \approx 351\MeV$,
close to the lepton triple's $v_0^2 \approx 314\MeV$. A third 
tuple $(c,b,t)$ stands apart, having no meson counterpart because 
top diquarks do not hadronize. We discuss the diquark sector---which 
fills the symmetric $\mathbf{15} + \overline{\mathbf{15}}$ of $SU(5)$ 
flavor, in conflict with Pauli statistics for point-like diquarks, 
proving that the diquark cannot be a $qq$ composite---and outline 
how SUSY breaking proceeds through two mechanisms---F-term (chiral
condensate, with Volkov--Akulov as its low-energy realization) and
D-term (electromagnetic Fayet--Iliopoulos)---plus a separate
heavy-sector origin for the $(c,b,t)$ triple. The O'Raifeartaigh
superpotential produces the Koide seed spectrum exactly when
$gv/m = \sqrt{3}$, reducing the energy-balance condition to a
single constraint on the couplings. The bloom from seed to full
triple doubles the vacuum charge, yielding
$\sqrt{m_b} = 3\sqrt{m_s} + \sqrt{m_c}$ (matching PDG to
$0.1\sigma$), and all quark Koide triples obligatorily mix
up-type and down-type mass eigenstates.
We construct the explicit $\mathcal{N}=1$ Lagrangian---K\"ahler
potential, superpotential (Seiberg constraint plus Yukawa with
separate $H_u$, $H_d$ doublets),
$D$-terms, and soft breaking from a spurion with
$\langle F\rangle = f_\pi^2$---and show that the mixed up/down-type
Koide triples require $\tan\beta = 1$ for the mass-level condition
to equal a Yukawa-level condition.  A negative quartic
K\"ahler correction $c = -1/12$ stabilizes the pseudo-modulus at
$\langle X\rangle = \sqrt{3}\,\mu$ (the K\"ahler pole), while the
$v_0$-doubling originates in the K\"ahler sector via a bion
correction $|\!\sum s_k\sqrt{\hat m_k}|^2$ from
monopole-instantons on $\mathbb{R}^3\times S^1$; a
three-instanton superpotential $(\det M)^3/\Lambda^{18}$
generates a $\cos(9\delta)$ potential whose minima include the
observed lepton Koide phase to $33\,\text{ppm}$ ($0.9\sigma$
in $m_\tau$).
After these constraints, only two to three
free parameters remain ($m_u$, $m_d$, $m_s$), with $m_u$ and
$m_d$ connected to the Cabibbo angle via the Oakes relation.
The $SO(32)$ adjoint accommodates the sBootstrap representations;
$E_8 \times E_8$ is blocked because $E_8$ decompositions produce
only antisymmetric $\mathbf{10}$, never the symmetric $\mathbf{15}$
required for diquarks. A Diophantine system on
the scalar content uniquely selects $N=3$ generations.
\end{abstract}

\tableofcontents

%======================================================================
\section{Introduction}
\label{sec:intro}
%======================================================================

The sBootstrap \cite{Rivero2005,Rivero2012} proposes that the Standard 
Model's pseudoscalar mesons and leptons fill common $\mathcal{N}=1$ 
supermultiplets. Its foundational relation is
\begin{equation}
\label{eq:pimu}
m_\pi^2 - m_\mu^2 \approx f_\pi^2\,,
\end{equation}
identifying the pion decay constant $f_\pi \approx 92.2\MeV$ as the 
SUSY-breaking order parameter. Numerically, $m_\pi^2 - m_\mu^2 = 
8{,}316\MeV^2$ versus $f_\pi^2 = 8{,}501\MeV^2$, a $2.2\%$ 
discrepancy---comparable to the precision of the quark Koide triples 
discussed below. In this paper we develop two themes 
that change the structure of the program.

First, \textbf{mass is the square of a signed charge}: $m = \zch^2$ 
where $\zch = z_0 + z_i$ can be negative. The Koide ratio $Q = 2/3$ 
is then not numerology but the energy-balance condition 
$\langle z_i^2\rangle = z_0^2$. With signed charges, pseudoscalar 
mesons satisfy Koide, and the spectrum organizes into \emph{seed 
triples} with a massless member that grow into full triples upon 
symmetry breaking.

Second, \textbf{the full sBootstrap has three Koide tuples, not one}, 
and the breaking structure must generate all three: a light sector 
$(0,s,c) \to (-s,c,b)$ seen in both quarks and mesons, a lepton 
sector $(e\!\approx\!0, \mu, \tau)$, and a heavy sector $(c,b,t)$ 
that decouples from meson physics because the top quark does not 
hadronize. The diquark partners of quarks, needed to complete the 
SUSY multiplets, fill the symmetric $\mathbf{15} + \overline{\mathbf{15}}$ 
of $SU(5)$ flavor---the representation produced naturally by the 
SUSY generator and by the $SO(32)$ decomposition. A kinematic 
theorem shows that no $J=0$, color-$\bar{3}$ diquark can have 
symmetric flavor---regardless of orbital angular momentum---so the 
sBootstrap diquark cannot be a $qq$ composite but must be an
algebraically or string-theoretically generated object.

Third, \textbf{the Koide condition constrains the superpotential}.
The O'Raifeartaigh model---the simplest realization of $F$-term
SUSY breaking with three chiral superfields---produces a fermionic
mass spectrum that is an exact Koide seed triple when the ratio of
the cubic coupling to the bilinear mass parameter satisfies
$gv/m = \sqrt{3}$. This single algebraic constraint is the
microscopic content of the energy-balance condition. Separately,
a Diophantine system on the scalar content ($rs = 2N$,
$r(r+1)/2 = 2N$, $r^2+s^2-1=4N$) uniquely selects $N=3$
generations, and a Casimir eigenvalue equation predicts
$\sin^2\theta_W = 0.22310$ to $0.4\sigma$.

\paragraph{New results in this paper.}
Compared to the prior publication \cite{Rivero2005}, the
following are first reported here: the signed-charge formalism
and the meson Koide triple $(-\pi, D_s, B)$; the seed-triple
concept and the universal ratio $2 - \sqrt{3}$; the
$v_0$-doubling relation $\sqrt{m_b} = 3\sqrt{m_s} + \sqrt{m_c}$;
the O'Raifeartaigh derivation $gv/m = \sqrt{3}$ and its
K\"ahler stabilization at the pole $c = -1/12$; the Pauli theorem
for symmetric diquarks; the bloom instability under perturbative
$R$-breaking; the dual Koide on $(d,s,b)$ and its connection to
the Seiberg seesaw; the bion K\"ahler mechanism for the
$v_0$-doubling; the self-consistent prediction chain
$m_s \to m_c \to m_b \to m_t$; and the scale relation
$M_0 \approx \fpi^2/27$ tying the lepton Koide scale
to the pion decay constant; and the prediction $\tan\beta = 1$
from the mixed up/down-type Koide triples.
The Casimir eigenvalue equation and the $SO(32)$
decomposition appeared in our prior work; here we add the
algebraic proof linking the Casimir ratio to the de~Vries form,
a uniqueness analysis of the equation, the
$N_f$-phase structure of the SQCD realization, and the
three-instanton origin of the $\cos(9\delta)$ potential
that selects the lepton Koide phase.

%======================================================================
\section{The Koide Formula as Energy Balance}
\label{sec:koide_algebra}
%======================================================================

\subsection{Charges, not masses}

The fundamental variable is the signed mass charge $\zch = z_0 + z_i 
\in \mathbb{R}$, with the physical mass $m = \zch^2 \geq 0$.

\subsection{Derivation of $Q = 2/3$}

For three charges $\zch_k = z_0 + z_k$ with $z_0$ the mean 
($\sum z_k = 0$):
\begin{equation}
Q = \frac{\sum \zch_k^2}{(\sum \zch_k)^2} 
= \frac{3z_0^2 + \sum z_k^2}{9z_0^2}
= \frac{1}{3} + \frac{\langle z_k^2\rangle}{3z_0^2}\,.
\end{equation}
Therefore:
\begin{equation}
\boxed{Q = \frac{2}{3} \quad\Longleftrightarrow\quad 
\langle z_k^2\rangle = z_0^2\,.}
\end{equation}

The two conditions are:
\begin{enumerate}
\item \textbf{Zero-sum perturbations:} $z_1 + z_2 + z_3 = 0$ 
(definition of $z_0$ as the mean charge).
\item \textbf{Energy balance:} The mean energy of the perturbations 
equals the energy of the unperturbed charge: 
$\langle z_k^2\rangle = z_0^2$.
\end{enumerate}

Condition~2 is the physics. It says the splitting scale matches the
base scale---the coefficient of variation equals unity:
$\sigma(\zch)/\bar{\zch} = 1$.
(Dimensionally, $\zch$ has units $\MeV^{1/2}$, so
$z_0^2$ has units MeV---a mass scale, not an energy density.
The ``energy balance'' terminology is metaphorical.) This is neither perturbative 
($\sigma \ll \bar{\zch}$, giving $Q \to 1/3$, all equal) nor 
dominant ($\sigma \gg \bar{\zch}$, giving $Q \to 1$, one mass 
dominates). The Koide value $2/3$ is the equipartition point.

\subsection{Why one mass can vanish}

If $\langle z_k^2\rangle = z_0^2$, a perturbation $z_k$ can reach 
$|z_k| \sim z_0$, making $\zch_k = z_0 + z_k \approx 0$. A vanishing 
mass is the \emph{natural extreme} of energy balance, not fine-tuning.

%======================================================================
\section{Seed Triples: The $(0, a, b)$ Structure}
\label{sec:seeds}
%======================================================================

\subsection{Analytical condition}

For a triple $(0, \zch_a, \zch_b)$ with one massless member, the 
Koide ratio becomes:
\begin{equation}
Q = \frac{\zch_a^2 + \zch_b^2}{(\zch_a + \zch_b)^2}\,.
\end{equation}
Setting $Q = 2/3$ and writing $r = \zch_a/\zch_b$ (with $r < 1$):
\begin{equation}
3(r^2 + 1) = 2(r + 1)^2 \quad\Longrightarrow\quad r^2 - 4r + 1 = 0\,,
\end{equation}
giving:
\begin{equation}
\boxed{r = 2 - \sqrt{3} \approx 0.26795\,.}
\end{equation}

Equivalently, in mass space: 
$m_{\text{light}}/m_{\text{heavy}} = (2 - \sqrt{3})^2 = 7 - 4\sqrt{3} 
\approx 0.07180$.

This is a \emph{universal} prediction: any Koide triple with one 
massless member has the same charge ratio between the two massive 
members, regardless of their absolute scale.

\subsection{The quark seed: $(0, s, c)$}

The strange and charm quark charges give:
\begin{equation}
\frac{\sqrt{m_s}}{\sqrt{m_c}} = \frac{9.67}{35.64} = 0.2712\,,
\qquad r_{\text{Koide}} = 0.2679\,, \qquad \text{deviation: } 1.2\%\,.
\end{equation}

This triple has $Q = 0.6644$, deviating from $2/3$ by only $0.35\%$.
Its Koide scale is $v_0^2 = (\zch_s + \zch_c)^2/9 \approx 228\MeV$.

\paragraph{Scale ambiguity in the quark seed.}
Throughout this paper, quark masses follow PDG conventions:
$m_s(2\GeV) = 93.4\MeV$, $m_c(m_c) = 1.27\GeV$,
$m_b(m_b) = 4.18\GeV$, $m_t = 172.76\GeV$ (pole mass). This is a
mixed convention, with the light quark evaluated at 2~GeV and heavier
quarks at their own scale. QCD mass ratios are RG invariants at
leading order (the anomalous dimension $\gamma_m$ is flavor-universal),
so $Q$ is scheme-independent. However, the seed ratio
$\sqrt{m_s}/\sqrt{m_c}$ uses masses at \emph{different} scales:
at a common $\overline{\text{MS}}$ scale the ratio shifts to
$\sim 0.29$, worsening the seed match from 1.2\% to $\sim 8\%$.
This is a significant limitation of the quark seed claim. The meson
seed, which uses unambiguous physical (pole) masses, is free of this
scheme issue and provides the cleaner test of the seed hypothesis.

A historical note: Harari, Haut, and Weyers \cite{Harari} observed 
a quark mass ratio pattern that, in the language of the present paper, 
corresponds to an approximate seed triple---$(0, d, s)$---at a time 
when the up quark mass was 
still thought to be possibly zero. With current masses, 
$\sqrt{m_d}/\sqrt{m_s} = 0.224$, deviating from $2 - \sqrt{3}$ by 
$16\%$---too large for a good Koide triple. But if $m_s$ were 
$\sim 150\MeV$ and $m_d \sim 7\MeV$ (values common in that era), the 
ratio improves considerably. The seed idea was already present, ahead 
of its time.

\subsection{The meson seed: $(0, \pi, D_s)$}

The pion and $D_s$ meson charges give:
\begin{equation}
\frac{\sqrt{m_\pi}}{\sqrt{m_{D_s}}} = \frac{11.81}{44.37} = 0.2663\,, 
\qquad r_{\text{Koide}} = 0.2679\,, \qquad \text{deviation: } 0.6\%\,.
\end{equation}

This gives $Q = 0.6679$, deviating by $0.18\%$. It is the 
\textbf{meson counterpart} of the quark seed $(0, s, c)$---both 
involve the strange-charm sector. Indeed, $D_s = c\bar{s}$ is the 
meson that directly contains the $s$ and $c$ quarks.

\subsection{The lepton triple as an ``almost-seed''}

The charged lepton triple $(e, \mu, \tau)$ has $\zch_e = 0.715\MeV^{1/2}$, 
which is small but nonzero. If $m_e$ were zero, we would need 
$\sqrt{m_\mu}/\sqrt{m_\tau} = 2 - \sqrt{3} = 0.2679$; the actual value 
is $0.2439$, deviating by $9\%$. So the leptons are \emph{not} an 
exact seed---the electron, while very light, is not massless. But the 
triple is clearly the same family: the near-critical Koide phase puts 
one member close to zero.

\subsection{Scale coincidence}

The $v_0^2$ values of the seed triples cluster remarkably:

\begin{center}
\begin{tabular}{lcc}
\toprule
Triple & $v_0^2$ (MeV) & $v_0$ ($\MeV^{1/2}$) \\
\midrule
$(e, \mu, \tau)$ & 314 & 17.72 \\
$(0, \pi, D_s)$ & 351 & 18.73 \\
$(0, s, c)$ & 228 & 15.10 \\
\bottomrule
\end{tabular}
\end{center}

The lepton and meson seed triples share a Koide scale around
310--350~MeV, squarely in the QCD scale range. This is the
sBootstrap connection made quantitative: \textbf{the lepton triple
and the meson seed triple operate at the same scale.}

For the lepton triple, the coincidence is numerically sharp:
\begin{equation}
M_0 \;\approx\; \frac{\fpi^2}{27}\,,
\end{equation}
to $0.3\%$: the right-hand side gives $314.9\MeV$ versus
the fitted $M_0 = 313.8\MeV$.  Equivalently,
$v_0 \approx \fpi/(3\sqrt{3})$.  If this identification is exact
and $\delta_0 \bmod 2\pi/3 = 2/9$ (the three-instanton value),
all three charged lepton masses follow from~$\fpi$ alone,
each predicted to~${\sim}0.3\%$.

%======================================================================
\section{A Tale of Two Tuples}
\label{sec:two_tuples}
%======================================================================

\subsection{From seed to full triple}

The seed triples grow into full triples upon symmetry breaking:

\begin{center}
\begin{tabular}{lccc}
\toprule
Sector & Seed triple & Full triple & $v_0^2$ (full) \\
\midrule
Quarks & $(0, s, c)$ & $(-s, c, b)$ & 913 MeV \\
Mesons & $(0, \pi, D_s)$ & $(-\pi, D_s, B)$ & 1230 MeV \\
Leptons & $(0 \approx e, \mu, \tau)$ & $(e, \mu, \tau)$ & 314 MeV \\
\bottomrule
\end{tabular}
\end{center}

In each case:
\begin{itemize}
\item The seed has a massless (or near-massless) member, two massive 
members in the ratio $2 - \sqrt{3}$, and $Q \approx 2/3$.
\item The full triple has \textbf{no} massless member: the zero has 
been lifted (to $m_e$, or the sign-flipped strange/pion mass). The 
$b$ quark and $B$ meson are always present as flavor degrees of 
freedom; what the breaking does is give them their \emph{distinctive 
masses}---the specific values that complete the Koide triple. Before 
breaking, the $b$-direction sits at the flat direction ($\zch = 0$); 
after, it acquires $\zch_b = 64.65\MeV^{1/2}$, the value selected 
by the energy balance with $c$ and the sign-flipped $s$. The bloom 
is not particle creation but \emph{mass generation}.
\end{itemize}

The ``tale of two tuples'' is this: \textbf{each sector presents
first as a seed with a massless member, then ``blooms'' into a
full triple when SUSY breaks.} The massless member is the necessary
initial condition; the breaking lifts it and generates the full spectrum.

\paragraph{The vacuum-charge doubling (new).}
The vacuum charge $v_0 = \tfrac{1}{3}\sum\zch_k$ of the quark seed
is $v_0^{(s)} = (\sqrt{m_s} + \sqrt{m_c})/3 = 15.10\MeV^{1/2}$.
For the full triple, $v_0^{(f)} = (-\sqrt{m_s} + \sqrt{m_c} +
\sqrt{m_b})/3 = 30.21\MeV^{1/2}$. Their ratio is
\begin{equation*}
  v_0^{(f)}/v_0^{(s)} = 2.0005 \approx 2\,.
\end{equation*}
If the bloom exactly doubles the vacuum charge, then
$\sqrt{m_b} = 3\sqrt{m_s} + \sqrt{m_c}$, giving
$m_b = 4{,}177\MeV$---matching the PDG value
$4{,}180 \pm 30\MeV$ to $0.1\sigma$ ($0.06\sigma$ including
propagated input uncertainties; see \S\ref{subsec:fterm}). This relation is
independent of the Koide condition: it constrains $m_b$ from
$(m_s, m_c)$ alone, and is more precise (0.07\%) than the
overlapping-triples prediction (0.5\%).  A tension exists, however:
the exact seed Koide gives $m_c = 1{,}301\MeV$ (not the PDG
$1{,}270$), and applying the $v_0$-doubling to this exact seed
value gives $m_b = 4{,}233\MeV$ ($1.8\sigma$ from PDG). The
$v_0$-doubling works \emph{with the physical} $m_c$, not the exact
seed value---so these two constraints are mutually inconsistent at the
stated precision. The resolution may be that the seed Koide fixes the
ratio $m_c/m_s$ only approximately (the 2.4\% charm deviation being a
higher-order correction to the $gv/m = \sqrt{3}$ condition), while the
$v_0$-doubling constrains $m_b$ directly from the physical masses. In
this reading, the seed condition is the approximate one and the
$v_0$-doubling the more fundamental. This hierarchy is not derived;
it is imposed by the data.

\subsection{Precision summary}

\begin{center}
\begin{tabular}{lcccc}
\toprule
Triple & $Q$ & Dev.\ from $2/3$ & $v_0^2$ (MeV) & $\delta$ (rad) \\
\midrule
\multicolumn{5}{c}{\textit{Seed triples}} \\
$(0, \pi, D_s)$ & 0.6679 & 0.18\% & 351 & $\approx 3\pi/4$ \\
$(0, s, c)$ & 0.6644 & 0.35\% & 228 & $\approx 3\pi/4$ \\
$(0 \approx e, \mu, \tau)$ & 0.6667 & 0.001\% & 314 & 2.317 \\
\midrule
\multicolumn{5}{c}{\textit{Full triples}} \\
$(e, \mu, \tau)$ & 0.66666 & 0.001\% & 314 & 2.317 \\
$(-\pi, D_s, B)$ & 0.6674 & 0.10\% & 1230 & 2.806 \\
$(-s, c, b)$ & 0.6750 & 1.24\% & 913 & 2.744 \\
$(c, b, t)$ & 0.6695 & 0.42\% & 29\,580 & 2.163 \\
\bottomrule
\end{tabular}
\end{center}

The Koide phases $\delta$ cluster in the range $2.2$--$2.8$~rad
($126^\circ$--$161^\circ$), all displaced from the seed value
$3\pi/4 \approx 2.356$. The lepton triple barely moved from its
seed; the quark and meson triples bloomed further. All phases
lie within one $\mathbb{Z}_3$ fundamental domain ($[0, 2\pi/3]$
modulo $2\pi/3$).

The lepton triple is unique: seed and full triple are nearly
identical, because the electron mass is small enough to serve as
the ``zero'' while being nonzero enough to make $Q$ extremely
precise. The $(-s,c,b)$ triple has the worst deviation (1.24\%);
this is expected because the bloom is a nonperturbative process
that does not exactly preserve $Q$. The $v_0$-doubling constrains
$m_b$ independently (to 0.07\% via $\sqrt{m_b} = 3\sqrt{m_s} +
\sqrt{m_c}$), so the 1.24\% Koide deviation is the residual after
the more fundamental $v_0$-constraint is satisfied.

\subsection{Orthogonality of seed and full triples}

Each Koide triple defines a perturbation vector 
$\vec{z} = (z_1, z_2, z_3)$ constrained to the plane $\sum z_k = 0$ 
with length $|\vec{z}|^2 = 3z_0^2$. The angle between seed and 
full triples measures how much the breaking rotates the spectrum 
in charge space:

\begin{center}
\begin{tabular}{lc}
\toprule
Pair of triples & Angle \\
\midrule
$(0,s,c)$ vs.\ $(-s,c,b)$ & $22^\circ$ \\
$(0,\pi,D_s)$ vs.\ $(-\pi,D_s,B)$ & $26^\circ$ \\
$(e,\mu,\tau)$ vs.\ $(0,\mu,\tau)$ & $0.8^\circ$ \\
$(0,s,c)$ vs.\ $(0,\pi,D_s)$ & $0.3^\circ$ \\
\bottomrule
\end{tabular}
\end{center}

Two results are notable. First, the quark seed $(0,s,c)$ and meson 
seed $(0,\pi,D_s)$ are \textbf{nearly parallel} ($0.3^\circ$): the 
meson seed is the composite image of the quark seed, with the GOR 
map preserving direction while changing scale. Second, the lepton 
triple is essentially its own seed ($0.8^\circ$)---the bloom has 
barely begun.


%======================================================================
\section{The Meson Koide Triple: $(-\pi, D_s, B)$}
\label{sec:mesonkoide}
%======================================================================

With the pion charge negative, the triple $(-\pi, D_s, B)$ satisfies:
\begin{equation}
Q(-\pi, D_s, B) = 
\frac{m_\pi + m_{D_s} + m_B}{(-\sqrt{m_\pi} + \sqrt{m_{D_s}} + \sqrt{m_B})^2}
= 0.6674\,,
\end{equation}
deviating from $2/3$ by $0.10\%$---the second most precise Koide 
triple known. The energy balance:
\begin{equation}
\frac{\langle z_k^2\rangle}{z_0^2} = 1.002\,.
\end{equation}

For comparison, the ``standard'' all-positive quark triples $(d,s,b)$ 
and $(u,c,t)$ deviate from $Q = 2/3$ by 9.7\% and 27\% respectively 
at $\overline{\text{MS}}$ masses. \textbf{The meson triple is more
precise than any quark triple}, including $(c,b,t)$ (0.42\%).

The pion's negative charge has physical meaning: as a pseudo-Goldstone
boson, the pion would be massless ($\zch = 0$) if chiral symmetry were
exact. It sits just below zero ($\zch_\pi = -11.81$), approaching zero
from the negative side as $m_q \to 0$. The sign of $\zch$ encodes
Goldstone nature: pseudo-Goldstone bosons sit on the negative side
of the critical point, having crossed through zero from positive
territory during chiral symmetry breaking.

\paragraph{Look-elsewhere correction.}
Eleven pseudoscalar mesons ($\pi$, $K$, $\eta$, $\eta'$, $D$,
$D_s$, $B$, $B_s$, $B_c$, $\eta_c$, $\eta_b$) yield
$\binom{11}{3} = 165$ triples, times
four independent sign patterns, giving 660 combinations.
Scanning all 660, only $(-\pi, D_s, B)$ and the closely related
$(-\pi, D_s, B_s)$ (deviation $0.14\%$) lie within $0.2\%$ of
$Q = 2/3$. A Monte Carlo with $10^7$ random triples drawn
from $[100,10{,}000]\MeV$ gives $p_{\text{single}} = 2.1 \times 10^{-4}$
for a match at $0.10\%$; correcting for 660 trials yields
$p_{\text{LEE}} = 0.13$ ($1.1\sigma$). The meson Koide triple
is therefore \emph{not} statistically significant on its own,
unlike the lepton triple ($2.8\sigma$ after LEE). Its interest
lies in the structural pattern: all six closest hits share the
sign assignment $(-\pi, D_{(s)}, B_{(s,c)})$, i.e.\ a negative
pion paired with charm-containing and bottom-containing mesons---the
combination predicted by the sBootstrap's $SU(5)$ flavor structure.

%======================================================================
\section{The Third Tuple: $(c, b, t)$}
\label{sec:cbt}
%======================================================================

\subsection{A tuple without mesons}

The triple $(c, b, t)$ with all positive signs gives $Q = 0.6695$, 
deviating by $0.42\%$. Its Koide scale $v_0^2 \approx 29.6\GeV$ 
is far above the QCD scale.

This triple has \textbf{no meson counterpart}: the top quark decays 
weakly ($t \to Wb$, lifetime $\sim 5 \times 10^{-25}$~s) before the 
strong interaction can bind it ($\tau_{\text{QCD}} \sim 1/\Lambda_{\text{QCD}} 
\sim 10^{-24}$~s). No $T$-meson ($t\bar{q}$) exists. The meson matrix 
$M^i{}_j$ extends only to $b$-flavored states; the top sector lives 
purely at the quark level.

\subsection{Diquarks and the top}

In the sBootstrap, diquarks $[q_i q_j]$ are the scalar partners of 
antiquarks $\bar{q}_k$, with the baryon acting as the SUSY generator 
\cite{CattoGursey, Miyazawa}. For five active flavors 
$(u, d, s, c, b)$, the diquarks form a representation of $SU(5)_f$.

The top diquarks $[tq]$ are kinematic objects only---they do not 
hadronize. Yet their masses should follow the same pattern as other 
diquarks: \textbf{there is no dynamical reason for top diquark masses 
to differ from the others}. The $(c,b,t)$ Koide triple governs the 
quark masses that enter the diquark sector, and its energy balance 
$\langle z_k^2\rangle = z_0^2$ constrains the possible diquark 
spectrum even though the top diquarks are never observed as bound 
states.

\subsection{The $(c,b,t)$ triple has no seed}

Unlike the $(0,s,c)$ and $(0,\pi,D_s)$ seeds, there is no obvious 
$(0, x, y)$ seed for the $(c,b,t)$ triple. No pair of quarks in this 
sector satisfies $\sqrt{m_{\text{light}}}/\sqrt{m_{\text{heavy}}} 
\approx 2 - \sqrt{3}$: one has $\sqrt{m_c}/\sqrt{m_t} = 0.086$ 
(too small) and $\sqrt{m_c}/\sqrt{m_b} = 0.551$ (too large).

This means $(c,b,t)$ is \emph{not} a ``bloomed seed'' in the same 
sense as $(-s,c,b)$. It may represent a fundamentally different 
structure---a triple that was never seeded by a massless member but 
instead emerged fully formed at a higher breaking scale, perhaps 
connected to electroweak symmetry breaking rather than to chiral 
symmetry breaking.

%======================================================================
\section{Mass Relations and Iterative Chains}
\label{sec:mass_chains}
%======================================================================

The preceding sections established the Koide energy-balance condition
$Q = 2/3$ and identified the triples that satisfy it. In this section
we quantify how special the lepton triple really is, introduce the
angular parametrization that reveals a hidden rational structure in
the Koide phase, and trace an iterative chain that connects the heavy
quarks to the strange sector. We report both successes and failures:
the scaling rule that generates the quark triple from the lepton
triple works once and then stops, and the iterative chain stalls at
a mass too light for any known fermion. These negative results
constrain the scope of the pattern.

\subsection{The lepton Koide formula: how special?}
Using PDG~2024 values $m_e = 0.51100\MeV$, $m_\mu = 105.658\MeV$,
$m_\tau = 1776.86 \pm 0.12\MeV$:
\begin{equation}
Q(e,\mu,\tau) = \frac{(\sqrt{m_e} + \sqrt{m_\mu} + \sqrt{m_\tau})^2}
{m_e + m_\mu + m_\tau} = 1.500014 \pm 0.000015\,,
\end{equation}
where $Q = (\sum\sqrt{m})^2 / \sum m$ is the reciprocal of the
convention in \S\ref{sec:koide_algebra}: the Koide condition is
$Q = 3/2$ here, equivalently $\sum m / (\sum\sqrt{m})^2 = 2/3$
there. We use the $3/2$ convention in this section to match
the most common form in the Koide literature.
The deviation from $3/2$ is $|Q - 3/2| = 1.38 \times 10^{-5}$,
corresponding to $0.91\sigma$ given the uncertainty dominated by
$\delta m_\tau$.

\paragraph{Look-elsewhere effect.}
The Standard Model contains 9~charged fermions (3~leptons, 6~quarks).
An exhaustive scan of all $\binom{9}{3} = 84$ mass triplets yields
the ranking in Table~\ref{tab:koide_scan}.

\begin{table}[ht]
\centering
\caption{Top-ranked SM triplets by proximity to $Q = 3/2$. The gap
between rank~1 and rank~2 is a factor of~450.}
\label{tab:koide_scan}
\begin{tabular}{clc}
\toprule
Rank & Triplet & $|Q - 3/2|$ \\
\midrule
1 & $(e, \mu, \tau)$ & $1.38 \times 10^{-5}$ \\
2 & $(c, b, t)$ & $6.2 \times 10^{-3}$ \\
\bottomrule
\end{tabular}
\end{table}

To assess the statistical weight of the lepton result, we performed a
Monte Carlo study: $10{,}000$ random mass sets of 9~fermions drawn
log-uniform on $[0.1, 200{,}000]\MeV$ were scanned for the best
triplet. The fraction with $|Q - 3/2|_{\min} \leq 1.38 \times 10^{-5}$
gives a look-elsewhere-corrected significance of $\mathbf{2.8\sigma}$.
This is suggestive but not conclusive: the lepton Koide formula is
unlikely to be a pure accident, but the evidence does not reach the
discovery threshold.

\subsection{The Koide angle parametrization}
Any triple satisfying $Q = 3/2$ exactly can be written
\cite{Koide1983}
\begin{equation}
\label{eq:koide_angle}
\sqrt{m_k} = \sqrt{M_0}\left(1 + \sqrt{2}\cos\!\left(
\frac{2\pi k}{3} + \delta_0\right)\right), \qquad k = 0,1,2\,,
\end{equation}
where $M_0 > 0$ is the scale parameter and $\delta_0 \in [0, 2\pi)$
is the Koide phase. This is not an approximation: any three
non-negative numbers whose $Q$-ratio equals $3/2$ are exactly
parametrized by some $(M_0, \delta_0)$.

Fitting to the charged lepton masses:
\begin{equation}
M_0 = 313.84\MeV\,, \qquad \delta_0 = 2.3166\;\text{rad}\,.
\end{equation}
The scale $M_0 \approx 314\MeV$ is the Koide scale $v_0^2$ of
\S\ref{sec:seeds}, now identified as the mass parameter of the
angular decomposition.

\paragraph{The phase is close to $2/9$.}
Unlike $Q = 2/3$, which is the algebraic consequence of the
energy-balance condition (\S\ref{sec:koide_algebra}), the phase
$\delta_0$ is a free parameter that the energy balance does
not fix. Its value is an empirical datum, not a derived quantity.
The parametrization~\eqref{eq:koide_angle} has a $2\pi/3$ periodicity
in $\delta_0$ (permuting the three masses). The physically
meaningful quantity is the reduced phase
\begin{equation}
\delta_0 \bmod \frac{2\pi}{3} = 0.22222\;\text{rad}\,.
\end{equation}
This is strikingly close to the fraction $2/9 = 0.22222\ldots$:
\begin{equation}
\delta_0 \bmod \frac{2\pi}{3} - \frac{2}{9} = +7.4 \times 10^{-6}\,,
\end{equation}
a residual of $\mathbf{33}$~\textbf{ppm} relative to $2/9$
(using PDG 2024 central masses; within the $m_\tau$ uncertainty
of $0.12\MeV$, the residual can vanish).

\paragraph{Statistical significance of the rational approximation.}
There are 128~``nice'' fractions $p/q$ with $q \leq 20$ and
$0 < p/q < 2\pi/3$. The probability that a random phase (uniform
on $[0, 2\pi/3)$) falls within $7.4 \times 10^{-6}$ of \emph{any}
of these 128 fractions gives a look-elsewhere-corrected significance
of $\mathbf{3.3\sigma}$. Evaluated specifically against the fraction
$2/9$ alone (i.e., without look-elsewhere correction), the
significance is $\mathbf{4.5\sigma}$.
The deviation from~$2/9$ is $0.9\sigma$ in the
$m_\tau$~uncertainty, consistent with an exact match.

\subsection{The scaling rule: $M_1 = 3M_0$, $\delta_1 = 3\delta_0$}
\label{subsec:scaling}

A natural question is whether other Koide triples can be obtained
from the lepton parameters by a simple scaling. The ansatz
\begin{equation}
M_1 = 3\,M_0 = 941.5\MeV\,, \qquad
\delta_1 = 3\,\delta_0 = 6.9497\;\text{rad}\,,
\end{equation}
inserted into~\eqref{eq:koide_angle}, predicts a quark triple.
Table~\ref{tab:scaling_pred} compares predictions to PDG values.

\begin{table}[ht]
\centering
\caption{Masses predicted by the $\times 3$ scaling rule vs.\ PDG
central values. The quark Koide triple uses a negative sign for
$\sqrt{m_s}$, as in \S\ref{sec:two_tuples}.}
\label{tab:scaling_pred}
\begin{tabular}{lccc}
\toprule
Quark & Predicted (MeV) & PDG (MeV) & Deviation \\
\midrule
$s$ & 92.3 & 93.4 & $1.2\%$ \\
$c$ & 1360 & 1270 & $7.1\%$ \\
$b$ & 4197 & 4180 & $0.4\%$ \\
\bottomrule
\end{tabular}
\end{table}

The agreement for $s$ and $b$ is at the percent level; charm shows a
larger ($7.1\%$) discrepancy. For the physical masses,
$M_0(-s,c,b)/M_0(e,\mu,\tau) = 2.91$, not exactly $3$---the
$\times 3$ rule is approximate. By contrast, the vacuum-charge
doubling (\S\ref{sec:two_tuples}) predicts $m_b$ to $0.1\sigma$.
The predicted triple satisfies the
Koide condition with a negative square root for the strange quark:
\begin{equation}
Q(-\sqrt{m_s},\;\sqrt{m_c},\;\sqrt{m_b})
= \frac{(-\sqrt{m_s} + \sqrt{m_c} + \sqrt{m_b})^2}{m_s + m_c + m_b}
= \frac{3}{2}\,,
\end{equation}
holding to machine precision (residual $2.2 \times 10^{-16}$), since
the masses are generated from the parametrization~\eqref{eq:koide_angle}
which enforces $Q = 3/2$ identically.

\paragraph{Renormalization caveat.}
The PDG quark masses used in the comparison are quoted at different
scales: $m_s(2\GeV)$, $m_c(m_c)$, $m_b(m_b)$, and $m_t$ as the pole
mass. The predicted masses emerge at a single internal scale set by
$M_1$. Running all masses to a common $\overline{\text{MS}}$ scale
would change the individual deviations in Table~\ref{tab:scaling_pred}
but not the structural point: the $\times 3$ scaling produces a
recognizable quark triple. As noted in \S\ref{subsec:exactseed},
QCD mass \emph{ratios} are RG invariants at leading order, so the
Koide condition itself is scheme-independent, but the comparison
of absolute masses to PDG values carries an $O(\alpha_s)$ ambiguity.

\subsection{The second scaling fails}
Iterating the scaling rule:
\begin{equation}
M_2 = 9\,M_0 = 2825\MeV\,, \qquad
\delta_2 = 9\,\delta_0 = 20.849\;\text{rad}\,,
\end{equation}
gives predicted masses $(478,\; 92.3,\; 16{,}377)\MeV$. Only the
middle value echoes $m_s$. The lightest ($478\MeV$) does not match
any known quark mass, and the heaviest ($16.4\GeV$) falls between
$m_b$ and $m_t$ with no candidate. No SM triplet accommodates these
values.

\textbf{The scaling stops at one step.} This is a negative result
that constrains any attempt to build a recursive tower of Koide
triples: the multiplicative rule $M \to 3M$, $\delta \to 3\delta$
is not a general principle but a one-off relation between the lepton
and quark sectors. Whether this single success is accidental or
reflects a deeper connection (perhaps related to the three colors
of QCD, or to the three generations) remains open.

\subsection{Iterative chain: descending from $t$ and $b$}
\label{subsec:chain}

An alternative approach to connecting the quark masses is purely
algebraic: given two known masses, solve the Koide equation
$Q = 3/2$ for the third. Starting from the top and bottom quarks:

\paragraph{Step~1: $Q(m_b, m_?, m_t) = 3/2$.}
The Koide equation with $(m_b, m_t)$ known is quadratic in
$\sqrt{m_?}$, yielding four branches (two signs for $\sqrt{m_?}$,
two roots of each quadratic). The solutions are:
\begin{equation}
m_? = \{1357,\; 60{,}280,\; 1.34 \times 10^6,\;
3.55 \times 10^6\}\MeV\,.
\end{equation}
Only the first branch gives a physically viable mass:
$m_? = 1357\MeV$, which we identify with the charm quark
(PDG: $m_c(m_c) = 1270\MeV$, deviation $6.8\%$).

\paragraph{Step~2: $Q(m_c, m_?, m_b) = 3/2$.}
With $m_c = 1357\MeV$ (the chain value, not the PDG value) and
$m_b = 4180\MeV$:
\begin{equation}
m_? = 92.17\MeV\,,
\end{equation}
which we identify with the strange quark. The negative square root
emerges naturally from the algebra---the solution selects the branch
with $\sqrt{m_?} < 0$ in the signed-charge framework, consistent
with the $(-s, c, b)$ structure of \S\ref{sec:two_tuples}.

\paragraph{Step~3: the chain stalls.}
Continuing from $(m_s, m_c)$ with $m_s = 92.17\MeV$:
\begin{equation}
m_? = 0.035\MeV\,.
\end{equation}
This is far too light for either $u$ ($2.16\MeV$) or $d$
($4.67\MeV$). The chain does not reach the first-generation quarks;
it stalls at a mass two orders of magnitude below $m_u$.

\paragraph{Step~4: cyclic echo.}
The stalled mass provides an algebraic consistency check.
Taking $m_{\text{stall}} = 0.035\MeV$ and $m_s = 92.17\MeV$, and
solving $Q(m_{\text{stall}}, m_?, m_s) = 3/2$:
\begin{equation}
m_? = 1357\MeV \quad \text{(to machine epsilon)}\,.
\end{equation}
The chain returns to its starting point: $t \to c \to s \to
\text{stall} \to c \to \cdots$. This is not a numerical coincidence
but an algebraic necessity: the Koide condition, combined with the
specific masses of $b$ and $t$, defines a closed orbit in
$(m_1, m_2, m_3)$ space. The cycle is:
\begin{multline*}
\cdots \;\to\; (b,t) \xrightarrow{Q=3/2} c
\;\to\; (c,b) \xrightarrow{Q=3/2} s \\
\;\to\; (s,c) \xrightarrow{Q=3/2} \text{stall}
\;\to\; (\text{stall}, s) \xrightarrow{Q=3/2} c
\;\to\; \cdots
\end{multline*}

The chain captures $c$, $b$, $s$ but misses $u$, $d$, and the
leptons entirely. The first generation is outside the algebraic
orbit seeded by the third generation.

\subsection{Cross-checks and master comparison}
Two independent routes lead to the quark masses: the
lepton-seeded scaling rule (\S\ref{subsec:scaling}) and the
top-bottom iterative chain (\S\ref{subsec:chain}).
Table~\ref{tab:master} compares all predictions.

\begin{table}[ht]
\centering
\caption{Master comparison of predicted quark masses (MeV). Chain~A
descends from $(t,b)$ via the Koide equation. Chain~B applies
the $\times 3$ scaling to the lepton triple. The small discrepancy
between the two ($1357$ vs.\ $1360\MeV$ for charm) arises from
different inputs: Chain~A uses the PDG $m_t$ and $m_b$ directly,
while Chain~B uses the fitted $(M_0, \delta_0)$ from leptons.}
\label{tab:master}
\begin{tabular}{lcccc}
\toprule
Quark & Chain A & Chain B & PDG & Status \\
 & (from $t,b$) & (from leptons) & & \\
\midrule
$s$ & 92.2 & 92.3 & 93.4 & $1\text{--}1.2\%$ deviation \\
$c$ & 1357 & 1360 & 1270 & $7\text{--}7.1\%$ deviation \\
$b$ & (input) & 4197 & 4180 & $0.4\%$ deviation \\
$t$ & (input) & --- & 172\,760 & --- \\
stall & 0.035 & --- & --- & no SM match \\
$M_2$ set & --- & $(478,\,92,\,16377)$ & --- & no SM match \\
\bottomrule
\end{tabular}
\end{table}

The convergence of the two chains on $m_s \approx 92\MeV$ and
$m_c \approx 1357\text{--}1360\MeV$ is the central positive result.
Both routes, starting from completely different inputs (lepton masses
vs.\ heavy quark masses), produce the same strange and charm masses
to within $0.2\%$ of each other. Against PDG values, the strange
quark prediction is excellent ($\sim 1\%$), and the bottom quark
prediction is even better ($0.4\%$); the charm prediction shows a
persistent $\sim 7\%$ overshoot.

The negative results are equally informative:
\begin{itemize}
\item The $\times 3$ scaling does not iterate: the $\times 9$ step
produces no viable triple.
\item The iterative chain does not reach the first generation: it
stalls at $0.035\MeV$, two orders of magnitude below $m_u$.
\item The chain is algebraically cyclic: it loops through
$(c, s, \text{stall})$ indefinitely, never escaping to new masses.
\end{itemize}
These failures delineate the boundary of the Koide pattern: it
connects the second and third quark generations (and, via the
scaling rule, the leptons), but does not extend to the first
generation or beyond the Standard Model spectrum. Whatever
dynamics underlies the energy-balance condition
$\langle z_k^2\rangle = z_0^2$ operates within a limited domain.

%======================================================================
\section{Composite Scalar Counting and Generation Uniqueness}
\label{sec:counting}
%======================================================================

The sBootstrap rests on a single postulate: every scalar of the
Supersymmetric Standard Model is a composite of two quarks (diquark,
color triplet) or of a quark--antiquark pair (meson, color singlet),
with the electric charge given by the sum of the constituents' charges.
We now show that matching the number of such composites to the
fermionic degrees of freedom of $N$ generations uniquely forces
$N = 3$.

\subsection{Setup and Diophantine system}
Let there be $N$ generations, each containing quarks of two types:
\emph{down} (charge $-1/3$) and \emph{up} (charge $+2/3$).  Not
every quark need participate as a string endpoint; we allow
$r$ light down-type quarks and $s$ light up-type quarks that form
composites.  The Standard Model with right-handed neutrinos has
$4N$ real (Weyl) fermionic degrees of freedom per charge sector in
each color channel.  Matching composites to fermions, charge sector
by charge sector, yields the following system.

\begin{enumerate}
\item \textbf{Charge $+1/3$ composites} (one up $\times$ one down,
analogous to charged sleptons): $rs$ complex scalars must match
$2N$ real fermionic d.o.f., giving
\begin{equation}
\label{eq:rs}
rs = 2N\,.
\end{equation}

\item \textbf{Charge $-2/3$ composites} (symmetric down--down pairs,
color sextet): $\binom{r+1}{2} = r(r+1)/2$ complex scalars must
match $2N$ real fermionic d.o.f., giving
\begin{equation}
\label{eq:rr}
\frac{r(r+1)}{2} = 2N\,.
\end{equation}

\item \textbf{Neutral composites} (classified by $SU(r+s)$, minus
the overall singlet trace; includes the partner of the right-handed
neutrino): $(r+s)^2 - 1 = r^2 + s^2 + 2rs - 1$ states, of which
$2rs$ are charged and already counted.  The neutral count is therefore
$r^2 + s^2 - 1$, to be matched to $4N$ real d.o.f., giving
\begin{equation}
\label{eq:neutral}
r^2 + s^2 - 1 = 4N\,.
\end{equation}
\end{enumerate}

\noindent
The alternative neutral equation $r^2 + s^2 - 1 = 2N$ (left-handed
neutrinos only) produces no positive-integer solution when combined
with \eqref{eq:rs}--\eqref{eq:rr}; thus the sBootstrap \emph{requires}
right-handed neutrinos.

\subsection{The uniqueness theorem}

\medskip
\noindent\textbf{Theorem 1} (Generation uniqueness).
\textit{The system \eqref{eq:rs}--\eqref{eq:neutral} has a unique
solution over the positive integers:}
\begin{equation}
\boxed{(N,\,r,\,s) = (3,\,3,\,2)\,.}
\end{equation}

\medskip
\noindent\textit{Proof.}
Equating the right-hand sides of \eqref{eq:rs} and \eqref{eq:rr}
gives $rs = r(r+1)/2$, hence
\begin{equation}
\label{eq:s_from_r}
s = \frac{r+1}{2}\,,
\end{equation}
which requires $r$ odd.  Substituting \eqref{eq:rs} and
\eqref{eq:s_from_r} into \eqref{eq:neutral}:
\[
r^2 + \Bigl(\frac{r+1}{2}\Bigr)^{\!2} - 1 \;=\; 4\,\cdot\,\frac{r\,(r+1)/2}{2}
\;=\; r(r+1)\,.
\]
Expanding the left side,
\[
r^2 + \frac{r^2 + 2r + 1}{4} - 1 = r^2 + r\,,
\]
and clearing the denominator:
\[
4r^2 + r^2 + 2r + 1 - 4 = 4r^2 + 4r
\quad\Longrightarrow\quad
r^2 - 2r - 3 = 0\,,
\]
i.e.\ $(r-3)(r+1) = 0$.  Since $r > 0$, we obtain $r = 3$, whence
$s = 2$ and $N = rs/2 = 3$.
\hfill$\square$

\medskip
The solution has $r + s = 5$ light quarks acting as string endpoints,
furnishing the fundamental representation of an $SU(5)$ flavor group.
The top quark is the unique ``elephant''---too heavy to serve as a
constituent of the scalar composites---in precise agreement with the
Standard Model, where $m_t \gg \Lambda_{\text{QCD}}$ prevents
top-flavored hadrons from forming.

\subsection{Partial solutions and the hexagonal-number family}
Equations \eqref{eq:rs} and \eqref{eq:rr} alone, without the
neutral constraint, require $2N$ to be a hexagonal number
$h_k = k(2k-1)$:
\[
2N = 1,\; 6,\; 15,\; 28,\; 45,\; 66,\; \ldots
\]
Of these, only $2N = 1$ and $2N = 6$ have $r \leq 16$ (the
asymptotic-freedom bound $2n_f < 33$ limits $r + s \leq 16$).
The first is inadmissible ($N = 1/2$); the second gives $N = 3$.

\medskip
\noindent\textbf{Remark} (The next hexagonal solution).
The case $(N, r, s) = (14, 7, 4)$ with
$r + s = 11$ light flavors and $2N = 28$ total satisfies
\eqref{eq:rs}--\eqref{eq:rr} but fails \eqref{eq:neutral}:
$r^2 + s^2 - 1 = 64 \neq 4 \times 14 = 56$.
Asymptotic freedom ($b_0 > 0$ in the QCD beta function) would
also require $2n_f < 33$, i.e.\ $n_f \leq 16$; with $N = 14$
generations one would need $2N = 28$ flavors, which violates the
bound decisively.  The neutral equation and the asymptotic-freedom
bound independently single out $N = 3$.

\subsection{Identification with $SU(5)$ representations}
With $r = 3$ and $s = 2$, the $r + s = 5$ light quarks form the
fundamental $\mathbf{5}$ of $SU(5)_f$.  The composites organize
into three $SU(5)$ representations:

\paragraph{Mesons (quark--antiquark, color singlet):}
The $\mathbf{5} \otimes \bar{\mathbf{5}}$ decomposes as
\begin{equation}
\mathbf{5} \otimes \bar{\mathbf{5}}
= \mathbf{24} \oplus \mathbf{1}\,.
\end{equation}
The $\mathbf{24}$ carries the sleptons (charged and neutral) and
the meson matrix.  Under
$SU(5) \supset SU(3)_d \times SU(2)_u$:
\begin{equation}
\mathbf{24} = (\mathbf{1},\mathbf{1})
+ (\mathbf{1},\mathbf{3})
+ (\mathbf{3},\mathbf{2})
+ (\bar{\mathbf{3}},\mathbf{2})
+ (\mathbf{8},\mathbf{1})\,.
\end{equation}

\paragraph{Diquarks (quark--quark, color triplet):}
The symmetric product
\begin{equation}
\mathbf{5} \otimes_S \mathbf{5}
= \mathbf{15}
\end{equation}
decomposes as
\begin{equation}
\mathbf{15} = (\mathbf{1},\mathbf{3})
+ (\mathbf{3},\mathbf{2})
+ (\mathbf{6},\mathbf{1})\,,
\end{equation}
and its conjugate $\overline{\mathbf{15}}$ provides the
antidiquarks.

\medskip
\noindent\textbf{Proposition 1} (Symmetric product is forced).
\textit{The counting equation \eqref{eq:rr} selects the
symmetric product $\mathbf{15}$, not the antisymmetric
$\overline{\mathbf{10}}$.}

\medskip
\noindent\textit{Proof.}
The charge $-2/3$ composites are down--down pairs.
Under $SU(3)_d$, the symmetric product
$\mathbf{3} \otimes_S \mathbf{3} = \mathbf{6}$ has dimension
$r(r+1)/2 = 6$, matching equation~\eqref{eq:rr}.  The antisymmetric
product $\mathbf{3} \otimes_A \mathbf{3} = \bar{\mathbf{3}}$ has
dimension $r(r-1)/2 = 3 \neq 2N = 6$.  The symmetric color sextet,
and hence the $\mathbf{15}$, is the unique choice consistent
with the sBootstrap counting.
\hfill$\square$

\subsection{Charge census of $\mathbf{24} \oplus \mathbf{15}
\oplus \overline{\mathbf{15}}$}
Treating the $SU(5)$ simultaneously as a Georgi--Glashow
unification group \cite{Rivero2005}, one assigns standard electric
charges via $Q_{\text{em}} = T_3 - Y/6$ to every component.
Table~\ref{tab:census} displays the resulting charge census for
$\mathbf{24} \oplus \mathbf{15} \oplus \overline{\mathbf{15}}$.

\begin{table}[ht]
\centering
\begin{tabular}{lcccccc@{\qquad}c}
\toprule
& \multicolumn{6}{c}{Number of states at $|Q_{\text{em}}|$} & \\
\cmidrule(lr){2-7}
& $0$ & $\tfrac{1}{3}$ & $\tfrac{2}{3}$ & $1$
& $\tfrac{4}{3}$ & $2$ & Dim \\
\midrule
$\mathbf{24}$ & 10 & $2\!\times\!3$ & ---
& $2\!\times\!1$ & $2\!\times\!3$ & --- & 24 \\[2pt]
$\mathbf{15}$ & 1 & 3 & 9
& 1 & --- & 1 & 15 \\[2pt]
$\overline{\mathbf{15}}$ & 1 & 3 & 9
& 1 & --- & 1 & 15 \\[2pt]
\midrule
Total & 12 & $2\!\times\!6$ & $2\!\times\!9$
& $2\!\times\!2$ & $2\!\times\!3$ & $2\!\times\!1$ & 54 \\
\bottomrule
\end{tabular}
\caption{Charge census of
$\mathbf{24} \oplus \mathbf{15} \oplus \overline{\mathbf{15}}$
of $SU(5)$.
Notation ``$2\!\times\!k$'' means $k$ states at $+|Q_{\text{em}}|$ and $k$
at $-|Q_{\text{em}}|$; the $\mathbf{24}$ is self-conjugate.  In the
$\mathbf{15}$, the 9 states at $|Q_{\text{em}}|=2/3$ split as 3 from the
$(\mathbf{3},\mathbf{2})_{1/6}$ component at $Q_{\text{em}} = +2/3$ and 6
from the $(\mathbf{6},\mathbf{1})_{-2/3}$ color sextet at
$Q_{\text{em}} = -2/3$; the $\overline{\mathbf{15}}$ is the conjugate.
Charges $\pm 4/3$ in the $\mathbf{24}$ arise from the leptoquark
components $(\mathbf{3},\mathbf{2})_{-5/6}$ and
$(\bar{\mathbf{3}},\mathbf{2})_{5/6}$.
The 12 neutral states comprise $8 + 1 + 1 = 10$ from the
$\mathbf{24}$ plus one each from $\mathbf{15}$ and
$\overline{\mathbf{15}}$.}
\label{tab:census}
\end{table}

\noindent
Three cross-checks tie the census to the Diophantine system:
\begin{itemize}
\item \textbf{Neutral composites:} The 12 states at $Q_{\text{em}} = 0$
match $r^2 + s^2 - 1 = 9 + 4 - 1 = 12 = 4N$.
\item \textbf{Charge $+1/3$ composites} (one up $\times$ one
down): $rs = 3 \times 2 = 6$ composites, matching equation
\eqref{eq:rs} with $rs = 6 = 2N$.  These states distribute as
3 from $(\bar{\mathbf{3}},\mathbf{2})_{5/6} \subset \mathbf{24}$
and 3 from $\overline{\mathbf{15}}$.
\item \textbf{Charge $-2/3$ composites} (symmetric down--down):
$r(r+1)/2 = 6$ symmetric pairs fill the color sextet
$(\mathbf{6},\mathbf{1})_{-2/3} \subset \mathbf{15}$, matching
equation~\eqref{eq:rr} with $r(r+1)/2 = 6 = 2N$.
\end{itemize}

\subsection{Anomaly cancellation and asymptotic freedom}
The representation $\mathbf{24} \oplus \mathbf{15} \oplus
\overline{\mathbf{15}}$ must be anomaly-free if it is to appear
in a consistent gauge theory.

\medskip
\noindent\textbf{Proposition 2} (Anomaly cancellation).
\textit{The cubic $SU(5)$ anomaly of $\mathbf{24} \oplus
\mathbf{15} \oplus \overline{\mathbf{15}}$ vanishes.}

\medskip
\noindent\textit{Proof.}
The anomaly coefficient $A(\mathbf{R})$ for $SU(n)$ satisfies
$A(\text{adjoint}) = 0$ (the adjoint is real),
$A(\mathbf{15}) = n + 4 = 9$, and
$A(\overline{\mathbf{15}}) = -(n+4) = -9$.  Therefore
\[
A(\mathbf{24}) + A(\mathbf{15}) + A(\overline{\mathbf{15}})
= 0 + 9 + (-9) = 0\,.
\]
\hfill$\square$

\medskip
\noindent\textbf{Proposition 3} (Asymptotic freedom).
\textit{An $SU(5)$ gauge theory with matter in the
$\mathbf{24} \oplus \mathbf{15} \oplus \overline{\mathbf{15}}$
is asymptotically free.}

\medskip
\noindent\textit{Proof.}
The one-loop beta-function coefficient is
\[
b_0 = \frac{11}{3}\,C_2(G) - \frac{2}{3}\,\sum_{\text{Weyl}} T(\mathbf{R})\,,
\]
where $C_2(G) = n = 5$ for $SU(5)$, $T(\mathbf{adj}) = n = 5$,
and $T(\mathbf{15}) = (n+2)/2 = 7/2$.  The Weyl fermion content
consists of one adjoint and one $\mathbf{15} + \overline{\mathbf{15}}$
(treating each as a Weyl fermion), giving
$\sum T = 5 + 7/2 + 7/2 = 12$.  Therefore
\[
b_0 = \frac{11}{3} \times 5 - \frac{2}{3} \times 12
= \frac{55 - 24}{3} = \frac{31}{3} > 0\,.
\]
\hfill$\square$

\medskip
Both properties are nontrivial: anomaly cancellation ensures
the gauge symmetry is quantum-mechanically consistent, and
$b_0 > 0$ ensures the theory is well-defined in the ultraviolet.
Together they confirm that the representation forced by the
sBootstrap counting is physically viable.

\subsection{Summary of the counting argument}
The logical chain is:
\begin{enumerate}
\item The sBootstrap postulate (composites $=$ scalars of the SSM)
yields three Diophantine equations in $(N, r, s)$.
\item The unique positive-integer solution is $(3, 3, 2)$:
three generations, with the symmetric $\mathbf{15}$ forced over
the antisymmetric $\overline{\mathbf{10}}$.
\item The solution organizes into $SU(5)$ representations
$\mathbf{24} \oplus \mathbf{15} \oplus \overline{\mathbf{15}}$,
which are anomaly-free and asymptotically free.
\item Right-handed neutrinos are required: without them, no
integer solution exists.
\item The top quark is excluded from the composite sector by
the counting itself ($s = 2 < N = 3$), matching its
kinematic decoupling ($m_t \gg \Lambda_{\text{QCD}}$).
\end{enumerate}

%======================================================================
\section{The Diquark Sector and SU(5) Flavor}
\label{sec:diquarks}
%======================================================================

\subsection{Counting partners}

In the full sBootstrap \cite{Rivero2005}, the scalar partners of 
the Standard Model fermions include both mesons ($q\bar{q}$) and 
diquarks ($qq$). For five active flavors, the meson matrix gives 
$5 \times 5 = 25$ scalars. The diquarks $[q_i q_j]$ form a 
representation of $SU(5)_f$:
\begin{equation}
\mathbf{5} \otimes \mathbf{5} = \overline{\mathbf{10}}_A 
\oplus \mathbf{15}_S\,.
\end{equation}

The symmetric $\mathbf{15}_S$ gives 15 diquarks including 
same-flavor states $[uu]$, $[dd]$, \ldots, while the antisymmetric 
$\overline{\mathbf{10}}_A$ gives 10 diquarks with $i \neq j$ only.

\subsection{The symmetric representation}
\label{subsec:sym15}

The sBootstrap SUSY generator acts asymmetrically on the two
fermion sectors. When it acts on a \emph{lepton}, it produces a
meson: $\mathcal{Q} \cdot \ell \sim q\bar{q}$. Since the quark and antiquark 
are distinguishable, the meson lies in the full 
$\mathbf{5} \otimes \bar{\mathbf{5}} = \mathbf{24} + \mathbf{1}$. 
There is no Pauli constraint, and the meson spectrum is perfectly 
reproduced.

When it acts on a \emph{quark}, it produces a diquark:
$\mathcal{Q} \cdot q_i \sim q_i q_j$. The bosonic SUSY generator applied to 
a single flavor index $i$ naturally produces the \textbf{symmetric} 
product---the $\mathbf{15}_S$ representation, not the antisymmetric 
$\overline{\mathbf{10}}$. This is the representation identified 
in the $SO(32)$ decomposition \cite{Rivero2005}: the adjoint of 
$SO(32)$ under $SU(5)_f$ decomposes to include the 
$\mathbf{15} + \overline{\mathbf{15}}$ (diquarks) alongside the 
$\mathbf{1} + \mathbf{24}$ (mesons).

The 15 includes 5 same-flavor states ($[uu]$, $[dd]$, $[ss]$, 
$[cc]$, $[bb]$) absent from the $\overline{\mathbf{10}}$. Among 
these are exotic $+4/3$-charged diquarks such as $[uu]$. The cost: 
an additional chiral $+4/3$ quark per generation, as identified 
in \cite{Rivero2005}. Whether these exotic states appear in nature 
as resonances, or whether they decouple at some high scale, is an 
open question.

\subsection{The Pauli theorem for scalar diquarks}
\label{subsec:pauli}

The requirement of the symmetric $\mathbf{15}$ conflicts with 
Fermi statistics---not merely for $L=0$ diquarks, but for 
\emph{all} scalar diquarks in the color $\bar{3}$ channel, 
regardless of internal orbital structure.

\medskip
\noindent\textbf{Theorem.}
\textit{For a spin-0 ($J=0$) diquark in the color $\bar{3}$ 
(antisymmetric) representation, the flavor wave function is 
antisymmetric for \textbf{all} values of the orbital angular 
momentum $L$ and spin $S$.}

\medskip
\noindent\textit{Proof.}
$J = 0$ requires $L = S$ (since $|L - S| \leq J \leq L + S$). 
The exchange symmetry of the spin-orbital wave function is 
$(-1)^{L+S+1} = (-1)^{2L+1} = -1$ for all $L$. Combined with 
the color $\bar{3}$ factor $(-1)$, the non-flavor product is 
$(-1)(-1) = +1$ (symmetric). Total antisymmetry then requires 
flavor to be antisymmetric: $\overline{\mathbf{10}}$. \hfill$\square$

\medskip
This result is purely kinematic---independent of dynamics, 
binding mechanisms, or spatial extent. No wave function, however 
extended, can produce a symmetric-flavor scalar diquark in color 
$\bar{3}$ while respecting Fermi statistics. The key cases:

\begin{center}
\begin{tabular}{ccccc}
\toprule
$L$ & $S$ & Spin-orbital $(-1)^{2L+1}$ & $\times$ Color & Flavor \\
\midrule
0 & 0 & $-1$ & $\times\,(-1) = +1$ & must be A \\
1 & 1 & $-1$ & $\times\,(-1) = +1$ & must be A \\
2 & 2 & $-1$ & $\times\,(-1) = +1$ & must be A \\
\bottomrule
\end{tabular}
\end{center}

The $L=1$, $S=1$ row closes the ``P-wave diquark'' escape route 
that might otherwise seem available.

The theorem has a sharp consequence: \textbf{the sBootstrap 
diquark in the $\mathbf{15}$ cannot be a $qq$ composite.} It 
must be either:
\begin{enumerate}
\item An algebraically generated object---the image of the 
Miyazawa supercharge \cite{Miyazawa, CattoGursey} acting on 
an antiquark, whose symmetry is determined by the superalgebra 
rather than by the statistics of constituent quarks; or
\item A fundamental string-theoretic object, as in the $SO(32)$ 
decomposition \cite{Rivero2005}, where the $\mathbf{15}$ 
arises from the adjoint representation of the gauge group and 
the open-string diquark is not built from point-like quarks.
\end{enumerate}

The contrast with the meson sector is instructive. When the 
SUSY generator acts on a lepton, it produces a meson 
($q\bar{q}$): the quark and antiquark are distinguishable, 
no Pauli constraint applies, and the meson spectrum is 
reproduced without tension. When it acts on a quark, it 
produces a diquark in the ``wrong'' representation---and the 
theorem shows this cannot be fixed by any internal structure 
of a $qq$ pair. The asymmetry between the two SUSY directions 
is absolute: mesons work as composites; diquarks do not. 
This is an unresolved obstruction: the sBootstrap requires
the symmetric $\mathbf{15}$, but QCD composite diquarks cannot
fill it. The proposed resolutions---Miyazawa supercharges and
SO(32) string embeddings---are directions for future work, not
established mechanisms. The diquark sector remains the weakest
point of the program.

%======================================================================
\section{Two Algebraic Structures}
\label{sec:two_structures}
%======================================================================

The sBootstrap particle content and mass spectrum are constrained by
two \emph{independent} algebraic structures.  Part~A
(\S\ref{subsec:so32_chain}--\ref{subsec:so32_open}) derives the
representation content from an $SO(32)$ embedding.  Part~B
(\S\ref{subsec:casimir_eq}--\ref{subsec:casimir_open}) derives an
electroweak mass ratio from a Casimir eigenvalue equation.
No algebraic connection between the two has been found
(see \S\ref{subsec:two_gap}).

%----------------------------------------------------------------------
\subsection{Part A: The SO(32) embedding chain}
\label{subsec:so32_chain}
%----------------------------------------------------------------------

The $SO(32)$ gauge group of Type~I / heterotic string theory
admits the maximal-subgroup chain
\begin{equation}
\label{eq:so32_chain}
SO(32) \;\supset\; SU(16) \times U(1)_A
  \;\supset\; SU(5) \times SU(3) \times U(1)_A \times U(1)_B \,.
\end{equation}
The first step is the standard embedding $SO(2n) \supset SU(n)
\times U(1)$; the second decomposes $SU(16)$ under the regular
subalgebra $SU(5) \times SU(3) \times U(1)_B$.
Here $SU(5)$ is the Georgi--Glashow unification group and $SU(3)$
is a \emph{horizontal} (flavor) symmetry acting on the three
light generations---it is \textbf{not} $SU(3)_c$.

The fundamental $\mathbf{32}$ of $SO(32)$ decomposes as
\begin{equation}
\label{eq:32decomp}
\mathbf{32} = \underbrace{(\mathbf{5},\mathbf{3})}_{R_1}
  \;\oplus\; \underbrace{(\bar{\mathbf{5}},\bar{\mathbf{3}})}_{R_2}
  \;\oplus\; \underbrace{(\mathbf{1},\mathbf{1})}_{R_3}
  \;\oplus\; \underbrace{(\mathbf{1},\mathbf{1})}_{R_4}\,,
\end{equation}
where $\dim = 15 + 15 + 1 + 1 = 32$ and the four summands carry
distinct $U(1)_A \times U(1)_B$ charges.

\subsection{Decomposition of the adjoint}
The adjoint of $SO(32)$ is the antisymmetric square of the
fundamental: $\mathbf{496} = \wedge^2(\mathbf{32})$.  Distributing
over the four summands in \eqref{eq:32decomp} and using the identity
\begin{equation}
\wedge^2(\mathbf{5},\mathbf{3})
  = \bigl[\wedge^2(\mathbf{5}) \otimes S^2(\mathbf{3})\bigr]
    \;\oplus\;
    \bigl[S^2(\mathbf{5}) \otimes \wedge^2(\mathbf{3})\bigr]
  = (\mathbf{10},\mathbf{6})
    \;\oplus\;
    (\mathbf{15},\bar{\mathbf{3}})\,,
\end{equation}
one obtains the full decomposition listed in
Table~\ref{tab:496}.

\begin{table}[ht]
\centering
\caption{Decomposition of the $SO(32)$ adjoint $\mathbf{496}$
under $SU(5) \times SU(3) \times U(1)_A \times U(1)_B$.}
\label{tab:496}
\begin{tabular}{lcrrc}
\toprule
$SU(5)\!\times\!SU(3)$ rep & Dim & $Q_A$ & $Q_B$ & Origin \\
\midrule
\multicolumn{5}{l}{\textit{Adjoint sector
  ($Q_A = Q_B = 0$, 226 states):}} \\[3pt]
$(\mathbf{24},\,\mathbf{8})$ & 192 & $0$ & $0$
  & $R_1 \otimes R_2$ \\
$(\mathbf{24},\,\mathbf{1})$ & 24  & $0$ & $0$
  & $R_1 \otimes R_2$ \\
$(\mathbf{1},\,\mathbf{8})$  & 8   & $0$ & $0$
  & $R_1 \otimes R_2$ \\
$(\mathbf{1},\,\mathbf{1})$  & 1   & $0$ & $0$
  & $R_1 \otimes R_2$ \\
$(\mathbf{1},\,\mathbf{1})$  & 1   & $0$ & $0$
  & $R_3 \otimes R_4$ \\
\midrule
\multicolumn{5}{l}{\textit{Matter sector
  ($Q_A,\,Q_B \neq 0$, 270 states):}} \\[3pt]
$(\mathbf{10},\,\mathbf{6})$  & 60 & $+2$ & $+2$
  & $\wedge^2(R_1)$ \\
$(\overline{\mathbf{10}},\,\bar{\mathbf{6}})$ & 60 & $-2$ & $-2$
  & $\wedge^2(R_2)$ \\
$(\mathbf{15},\,\bar{\mathbf{3}})$ & 45 & $+2$ & $+2$
  & $\wedge^2(R_1)$ \\
$(\overline{\mathbf{15}},\,\mathbf{3})$ & 45 & $-2$ & $-2$
  & $\wedge^2(R_2)$ \\
$(\mathbf{5},\,\mathbf{3})_{+}$  & 15 & $+2$ & $-14$
  & $R_1 \otimes R_3$ \\
$(\mathbf{5},\,\mathbf{3})_{-}$  & 15 & $0$  & $+16$
  & $R_1 \otimes R_4$ \\
$(\bar{\mathbf{5}},\,\bar{\mathbf{3}})_{+}$ & 15 & $-2$ & $+14$
  & $R_2 \otimes R_3$ \\
$(\bar{\mathbf{5}},\,\bar{\mathbf{3}})_{-}$ & 15 & $0$  & $-16$
  & $R_2 \otimes R_4$ \\
\midrule
\textbf{Total} & \textbf{496} & & & \\
\bottomrule
\end{tabular}
\end{table}

\subsection{Physical content and the projection problem}
\label{subsec:so32_open}

Restricting to the adjoint sector $Q_A = Q_B = 0$ and dropping the
horizontal $SU(3)$ factor, one reads off the $SU(5)$ content:
\begin{equation}
\mathbf{24} \;\oplus\; \mathbf{1}\,,
\end{equation}
which is precisely the meson/slepton content of the sBootstrap.
The matter sector at $|Q_A| = 2$ contains
\begin{equation}
\mathbf{15} \;\oplus\; \overline{\mathbf{15}}
  \;\oplus\; \mathbf{10} \;\oplus\; \overline{\mathbf{10}}
  \;\oplus\; \mathbf{5} \;\oplus\; \bar{\mathbf{5}}\,,
\end{equation}
among which the $\mathbf{15} + \overline{\mathbf{15}}$ are
the symmetric diquarks demanded by the sBootstrap counting
(\S\ref{subsec:sym15}).

Thus the sBootstrap representations appear in the $SO(32)$
adjoint: the $\mathbf{24}$ in the neutral sector ($Q_A\!=\!Q_B\!=\!0$)
and the $\mathbf{15} \oplus \overline{\mathbf{15}}$ in the matter
sector ($|Q_A|\!=\!2$). However, the remaining $496 - 54 = 442$
states (the $\mathbf{10}$, $\overline{\mathbf{10}}$,
$\mathbf{5}$, $\bar{\mathbf{5}}$, the SU(3)-horizontal copies,
and the neutral singlets) must be projected out or made massive.
The mechanism for this projection---whether by an orbifold,
a GSO-type condition, or a flux background---is an open problem.

This embedding was first noted in \cite{Rivero2005}.  Its role
here is to provide a UV origin for the symmetric $\mathbf{15}$
that the Pauli theorem (\S\ref{subsec:pauli}) forbids at the
level of point-particle composites: in the string picture,
diquarks are open-string states whose symmetry is determined by
the gauge group, not by Fermi statistics of constituents.

\paragraph{Why $SO(32)$ and not $E_8 \times E_8$.}
Anomaly cancellation in ten dimensions selects two gauge groups:
$SO(32)$ and $E_8 \times E_8$. The latter cannot accommodate the
sBootstrap. Every decomposition of $E_8$ under $SU(5)$ produces
the antisymmetric $\mathbf{10}$, never the symmetric $\mathbf{15}$:
the maximal-subgroup chain $E_8 \supset SU(5) \times SU(5)$ gives
$\mathbf{248} = (\mathbf{24},\mathbf{1}) + (\mathbf{1},\mathbf{24})
+ (\mathbf{10},\bar{\mathbf{5}}) + (\overline{\mathbf{10}},\mathbf{5})
+ (\mathbf{5},\mathbf{10}) + (\bar{\mathbf{5}},\overline{\mathbf{10}})$,
with no $\mathbf{15}$ appearing anywhere.  The obstruction is
irreducible: it reflects the fact that $E_8$'s simply-laced root
lattice admits only antisymmetric tensor representations of
$SU(5)$.  By contrast, $SO(32)$'s adjoint $\mathbf{496} =
\wedge^2(\mathbf{32})$ naturally contains the symmetric
$\mathbf{15}$ through the identity
$\wedge^2(\mathbf{5} \otimes \mathbf{3}) \supset
S^2(\mathbf{5}) \otimes \wedge^2(\mathbf{3})
= (\mathbf{15}, \bar{\mathbf{3}})$.  The Type~I realization, where
quark Chan-Paton factors directly generate $SO(32)$ from
$5\!\times\!3 + 1 = 16$ endpoint labels, is the most natural home.

The $E_8 \times E_8$ heterotic string is connected to Type~I $SO(32)$
by a chain of dualities (S-duality to heterotic $SO(32)$, then
T-duality on a non-simply-connected torus).  These dualities may
allow the sBootstrap content to appear in the $E_8 \times E_8$ frame
through a non-standard embedding, but the direct decomposition
route is blocked by the $\mathbf{15}$ vs.\ $\mathbf{10}$
obstruction.

%----------------------------------------------------------------------
\subsection{Part B: The Casimir eigenvalue equation}
\label{subsec:casimir_eq}
%----------------------------------------------------------------------

Consider the eigenvalue equation
\begin{equation}
\label{eq:casimir_quartic}
x^2 \;+\; C_2\,x \;-\; C_2 = 0\,,
\qquad C_2 = s(s+1)\,,
\end{equation}
where $C_2$ is the quadratic Casimir of $SU(2)$ at spin $s$ and
$x = M^2/m^2$ for a mass scale $m$ to be determined.  Equivalently,
$x^2/(1-x) = C_2$: the eigenvalue is a self-consistent solution
where the ``squared coupling'' $x^2$ equals $C_2$ times the
complement $(1-x)$.  The solutions are
\begin{equation}
\label{eq:casimir_roots}
x_{\pm}(s) = \frac{-C_2 \pm \sqrt{C_2^2 + 4C_2}}{2}\,.
\end{equation}

\paragraph{Structural constraints (codimension analysis).}
The most general two-parameter family of such equations is
\begin{equation}
x^2 + (f_0 + f_1 C_2)\,x + (g_0 + g_1 C_2) = 0\,,
\end{equation}
with four real coefficients $(f_0, f_1, g_0, g_1)$.  Three
structural conditions pin the polynomial to \eqref{eq:casimir_quartic}:
\begin{enumerate}
\item $f_0 = 0$: no spin-independent mass term;
\item $g_0 = 0$: a massless eigenstate exists at $s = 0$;
\item $g_1 = -f_1$: coefficient antisymmetry
  (the constant and linear terms in $C_2$ are related).
\end{enumerate}
Together with the normalization $f_1 = 1$, these exhaust the
freedom.  A systematic scan of 582 polynomial alternatives
(quadratics, cubics, quartics with Casimir-type functions and
rational coefficients) confirms uniqueness: no other
equation of this class reproduces $\sin^2\!\theta_W$ to
sub-percent accuracy at $(s_1, s_2) = (1/2, 1)$.
The ratio derived below is therefore a
\emph{consequence} of these constraints, not a free parameter.

\subsection{The electroweak ratio}
Define the ratio
\begin{equation}
R \;\equiv\; 1 - \frac{x_+(1/2)}{x_+(1)}\,.
\end{equation}
Evaluating \eqref{eq:casimir_roots} at $s = 1/2$ ($C_2 = 3/4$)
and $s = 1$ ($C_2 = 2$):
\begin{align}
x_+(1/2) &= \tfrac{1}{8}\bigl(-3 + \sqrt{57}\bigr)
  \,,\quad\text{so } x_+(1/2)\,m^2 = M_W^2\,,
\\[4pt]
x_+(1) &= -1 + \sqrt{3}
  \,,\quad\text{so } x_+(1)\,m^2 = M_Z^2\,.
\end{align}
The identification with $W$ and $Z$ is motivated by their
$SU(2)$ quantum numbers ($s = 1/2$ maps to the doublet structure of $W$,
$s = 1$ to the triplet).  Then
\begin{equation}
R = 1 - \frac{x_+(1/2)}{x_+(1)}
  = \frac{(\sqrt{19} - 3)(\sqrt{19} - \sqrt{3}\,)}{16}
  \;\approx\; 0.2231013\ldots
\end{equation}
The factored form follows from the discriminants of
\eqref{eq:casimir_quartic}: $57 = 3 \times 19$ at $s = 1/2$ and
$12 = 4 \times 3$ at $s = 1$.  Rationalizing the denominator
$\sqrt{3} - 1$ yields an expression in $\sqrt{19}$ and $\sqrt{3}$
alone, which coincides term-by-term with the form found
empirically by de~Vries.

This should be compared with the on-shell electroweak mixing
parameter:
\begin{equation}
\label{eq:sin2_onshell}
\sin^2\!\theta_W^{\text{os}}
  \;\equiv\; 1 - \frac{M_W^2}{M_Z^2}
  = 0.22320 \pm 0.00026\,,
\end{equation}
using $M_Z = 91.1876 \pm 0.0021\GeV$ and
$M_W = 80.369 \pm 0.013\GeV$ (PDG~2024).

Since $R = 1 - x_+(1/2)/x_+(1)$ and
$\sin^2\!\theta_W = 1 - M_W^2/M_Z^2$, the two expressions are
\emph{algebraically identical} once one identifies
$x_+(s) = M^2(s)/m^2$.  The prediction $R \approx
\sin^2\!\theta_W$ and the $W$-mass prediction derived below are
therefore \textbf{one and the same observation}, not two
independent tests.

\paragraph{$W$-mass prediction.}
From $M_W = M_Z\sqrt{1 - R}$:
\begin{equation}
M_W^{\text{pred}}
  = 91.1876 \times \sqrt{1 - 0.22310}
  = 80.374 \pm 0.002\GeV\,,
\end{equation}
which lies $0.39\sigma$ from the PDG average
$M_W = 80.369 \pm 0.013\GeV$ and strongly disfavors the
anomalous CDF-II measurement $M_W^{\text{CDF}} = 80.4335 \pm
0.0094\GeV$ (6.2$\sigma$ tension with $R$).

\paragraph{Scale parameter.}
Setting $x_+(1)\,m^2 = M_Z^2$ gives
\begin{equation}
m = \frac{M_Z}{\sqrt{\sqrt{3} - 1}}
  = 106.577\GeV
  \approx m_\mu \times 10^3\,,
\end{equation}
a coincidence whose significance is unclear.

\subsection{Open issues for the Casimir equation}
\label{subsec:casimir_open}

\paragraph{Renormalization scheme dependence.}
The comparison \eqref{eq:sin2_onshell} uses the on-shell
definition $\sin^2\!\theta_W = 1 - M_W^2/M_Z^2 = 0.22320$,
\textbf{not} the $\overline{\text{MS}}$ value
$\sin^2\!\theta_W(\overline{\text{MS}}) = 0.23122 \pm 0.00004$.
The two differ by $\Delta \approx 0.008$, far larger than
experimental errors.  The algebraic argument
\eqref{eq:casimir_quartic} contains no dynamical mechanism to
prefer one scheme over the other.  Until such a mechanism is
identified, the on-shell agreement should be regarded as
suggestive but scheme-dependent.

\paragraph{The fourth eigenvalue.}
At $s = 1/2$, the second root of \eqref{eq:casimir_quartic} gives
\begin{equation}
|M(1/2, -)|
  = m\,\sqrt{|x_-(1/2)|}
  = 122.4\GeV\,,
\end{equation}
which is 2.3\% below $M_H = 125.25 \pm 0.17\GeV$.  However,
this eigenstate carries spin $1/2$, while the Higgs boson is
spin $0$.  The near-degeneracy is noted as a numerical curiosity;
we do not claim a connection.

\paragraph{Origin.}
This observation was first made by H.~de~Vries in a Physics
Forums discussion (November 2004) and subsequently studied in
\cite{Rivero2005, Rivero2005b}.

%----------------------------------------------------------------------
\subsection{Independence of the two structures}
\label{subsec:two_gap}
%----------------------------------------------------------------------

Parts~A and~B describe independent algebraic structures that share
the same physical particles:
\begin{itemize}
\item The $SO(32)$ embedding (\S\ref{subsec:so32_chain}) governs
  the \emph{representation content}---which $SU(5)$ multiplets
  appear.
\item The Casimir equation (\S\ref{subsec:casimir_eq}) governs
  the \emph{mass/coupling pattern}---what numerical ratios emerge
  among the gauge-boson masses.
\item The two structures live over different number fields:
  the adjoint decomposition involves only rational
  dimensions, while the eigenvalue ratio $R$ lies in
  $\mathbb{Q}(\sqrt{3},\sqrt{19}\,)$---an extension not
  generated by any invariant of $SO(32)$.
\end{itemize}
No algebraic homomorphism linking the embedding chain
\eqref{eq:so32_chain} to the eigenvalue equation
\eqref{eq:casimir_quartic} has been found.  The gap between
\emph{content} (which states exist) and \emph{dynamics}
(what masses they acquire) remains the central open problem of
the sBootstrap program.

%======================================================================
\section{SUSY Breaking: From Seeds to Spectrum}
\label{sec:breaking}
%======================================================================

The SUSY-breaking story begins with the foundational observation 
\eqref{eq:pimu}: $m_\pi^2 - m_\mu^2 \approx f_\pi^2$. This is the 
Volkov--Akulov mass formula for a pseudo-Goldstino---the muon---whose 
superpartner---the pion---is a pseudo-Goldstone boson. The SUSY-breaking 
scale is $f_\pi \approx 92\MeV$, and the breaking mechanism is chiral: 
$\qqbar \neq 0$ breaks both chiral symmetry and the sBootstrap SUSY 
simultaneously. The Volkov--Akulov realization is the \emph{physical 
starting point}; what follows in this section is the SQCD machinery 
that makes it precise.

The dynamics must accomplish three things: generate the seed-to-bloom 
transition for each tuple, preserve the energy-balance condition 
$\langle z_k^2\rangle = z_0^2$ through the breaking, and produce 
the correct splitting between meson and lepton masses at scale $f_\pi$.

\subsection{The Seiberg framework}

The natural arena is $\mathcal{N}=1$ SQCD with gauge group $SU(3)_c$ 
and $N_f$ flavors of chiral superfields $Q^i$, $\bar{Q}_j$ in the 
fundamental and antifundamental representations. For $N_f = N_c = 3$ 
(the case relevant to light QCD), Seiberg's exact results 
\cite{Seiberg1995} establish that the theory confines, with the 
low-energy degrees of freedom being the composite meson superfield
\begin{equation}
M^i{}_j = Q^i \cdot \bar{Q}_j\,,
\end{equation}
a $3\times 3$ matrix of chiral superfields, together with baryons 
$B = \epsilon_{ijk}Q^iQ^jQ^k$ and $\bar{B}$. These satisfy the 
quantum-modified constraint
\begin{equation}
\label{eq:seiberg}
\det M - B\bar{B} = \Lambda^{6}\,,
\end{equation}
where $\Lambda$ is the dynamical scale ($2N_c = 6$ for $N_c = 3$).

Each entry $M^i{}_j$ is a chiral superfield containing a complex 
scalar (the meson) and a Weyl fermion (the lepton, in the sBootstrap 
identification). The scalar and fermion are degenerate in the SUSY 
limit; SUSY breaking splits them.

\subsection{The three chiral superfields of each seed}

The seed triple $(0, \zch_a, \zch_b)$ maps naturally onto three 
chiral superfields with the following mass structure before breaking:

\begin{center}
\begin{tabular}{lcl}
\toprule
Superfield & Mass & Role \\
\midrule
$\Phi_0$ & $m_0 = 0$ & Massless: flat direction \\
$\Phi_a$ & $m_a = \zch_a^2$ & Light massive member \\
$\Phi_b$ & $m_b = \zch_b^2$ & Heavy massive member \\
\bottomrule
\end{tabular}
\end{center}

The ratio $\zch_a/\zch_b = 2 - \sqrt{3}$ is enforced by the 
Koide energy-balance condition applied to the seed.

In the Seiberg description, these three superfields are entries of 
the meson matrix $M^i{}_j$. For the $(0, s, c)$ seed, the relevant 
entries involve the strange and charm quarks; $\Phi_0$ is the entry 
whose quark-mass combination vanishes in the seed limit (i.e., the 
direction in flavor space where $m_i + m_j \to 0$ through GOR).

\subsection{F-term breaking: the O'Raifeartaigh mechanism}
\label{subsec:fterm}

The VA relation \eqref{eq:pimu} tells us \emph{that} SUSY is broken at 
scale $f_\pi$. To understand \emph{how}, we need the microscopic 
mechanism. The O'Raifeartaigh model \cite{ORaifeartaigh} is the simplest 
realization of spontaneous $F$-term SUSY breaking. With three chiral 
superfields, the superpotential
\begin{equation}
W = f\,\Phi_0 + m\,\Phi_1\Phi_2 + g\,\Phi_0\Phi_1^2
\end{equation}
gives $F$-terms:
\begin{align}
F_{\Phi_0}^* &= f + g\,\phi_1^2\,, \\
F_{\Phi_1}^* &= m\,\phi_2 + 2g\,\phi_0\phi_1\,, \\
F_{\Phi_2}^* &= m\,\phi_1\,.
\end{align}
If $f \neq 0$ and $m \neq 0$, then $F_{\Phi_2}^* = 0$ requires 
$\phi_1 = 0$, but then $F_{\Phi_0}^* = f \neq 0$: SUSY is 
\textbf{necessarily broken}. The key features are:
\begin{enumerate}
\item A massless field ($\Phi_0$) exists at tree level---it is the 
pseudo-modulus.
\item The breaking is triggered because $F_{\Phi_0}$ cannot vanish.
\item After including loop corrections, $\Phi_0$ acquires a mass 
(the Coleman--Weinberg potential lifts the flat direction).
\end{enumerate}

The structural parallel with the seed triples is:

\begin{center}
\begin{tabular}{llll}
\toprule
O'Raifeartaigh & Flavor & Seed & After breaking \\
\midrule
$\Phi_0$ (massless) & $b$ / $B$ & $\zch = 0$ (flat direction) & $\zch_b = 64.65$ (mass generated) \\
$\Phi_1$ (light) & $s$ / $\pi$ & $\zch_s = 9.67$ & $\zch \to -9.67$ (sign-flipped) \\
$\Phi_2$ (heavy) & $c$ / $D_s$ & $\zch_c = 35.64$ & $\zch_c \approx 35.64$ (shifted) \\
\midrule
 & & $f = 0$ (SUSY) & $f = f_\pi$ (broken) \\
\bottomrule
\end{tabular}
\end{center}

This parallel is more than structural: the minimal O'Raifeartaigh
model gives $\Phi_1$ and $\Phi_2$ a single mass parameter $m$, but the
\emph{fermionic} mass matrix at the vacuum ($\phi_1 = 0$,
$\phi_0 = v$) is
\begin{equation*}
  W_{ij} = \begin{pmatrix} 0 & 0 & 0 \\ 0 & 2gv & m \\ 0 & m & 0 \end{pmatrix},
\end{equation*}
with eigenvalues $0$ (the goldstino) and $m_\pm = \sqrt{g^2v^2 + m^2} \pm gv$.
The mass ratio $m_-/m_+ = (\sqrt{t^2+1} - t)^2$, where $t = gv/m$,
is \emph{not} unity.  Setting $\sqrt{m_-}/\sqrt{m_+} = 2 - \sqrt{3}$
(the Koide seed ratio) requires $\sqrt{t^2+1} - t = 2 - \sqrt{3}$,
which gives
\begin{equation*}
  t = \frac{gv}{m} = \sqrt{3}\,,
\end{equation*}
the unique value at which $t^2 + 1 = 4$, so $m_\pm = (2 \pm \sqrt{3})\,m$.
The Koide ratio for the triple $(0, m_-, m_+) = (0,\, (2-\sqrt{3})\,m,
\, (2+\sqrt{3})\,m)$ evaluates to $Q = 4m/(\sqrt{2-\sqrt{3}} +
\sqrt{2+\sqrt{3}})^2 = 4/6 = 2/3$ exactly, since
$\sqrt{2-\sqrt{3}}\cdot\sqrt{2+\sqrt{3}} = 1$ and the squared sum
telescopes to $6$.

The Koide condition on the seed is therefore \emph{one algebraic
constraint on the superpotential parameters}: $gv = \sqrt{3}\,m$.
The coupling $g$, the pseudo-modulus VEV $v$, and the mass
parameter $m$ are individually undetermined, but their ratio is fixed.
This is the microscopic content of the energy-balance condition
$\langle z_k^2\rangle = z_0^2$ applied to the seed.

The $b$ quark and $B$ meson are
always present as flavor degrees of freedom; the seed configuration
$(0, s, c)$ means the $b$-direction sits at the flat direction before
breaking. The bloom is not particle creation but mass generation:
breaking lifts $\Phi_0$ from zero to $\zch_b \neq 0$, simultaneously
flipping the sign of $\zch_s$, so that the same three flavors
rearrange from $(0, s, c)$ into $(-s, c, b)$.

\paragraph{The bloom as a phase rotation.}
In the Koide parametrization~\eqref{eq:koide_angle}, the seed sits at
$\delta = 3\pi/4$, where $\cos(3\pi/4) = -1/\sqrt{2}$ forces one
mass to vanish: $\sqrt{m_0} = \sqrt{M_0}(1 + \sqrt{2}\cos\delta) = 0$.
A Koide-preserving bloom ($Q = 2/3$ maintained) is a pure rotation
$\delta: 3\pi/4 \to \delta_{\text{phys}}$ at fixed scale $v_0$
and fixed amplitude $r = \sqrt{2}\,v_0$. The total mass
$M = 6v_0^2$ is conserved---mass merely redistributes among the three
eigenvalues. The sign flip is instantaneous: the seed sits exactly at
the boundary $\zch_0 = 0$, and any infinitesimal rotation generates
a nonzero charge with definite sign.

Explicit R-symmetry breaking in the O'Raifeartaigh model (adding a
term $\delta W \sim \varepsilon\,\Phi_0\Phi_2$) lifts the goldstino
mass as $m_0 \sim \varepsilon^2/m$, but simultaneously pushes
$Q$ away from $2/3$---no nonzero $\varepsilon$ preserves the Koide
condition. The Koide seed is an \emph{unstable fixed point} of the
R-breaking perturbation. In the dynamical (SQCD) realization, the
bloom must therefore be a nonperturbative rearrangement---not a
smooth deformation of the O'Raifeartaigh spectrum.

A striking difference from standard O'Raifeartaigh: the 
pseudo-modulus $\Phi_0$ normally acquires the \emph{smallest} mass 
in the spectrum (a loop-suppressed Coleman--Weinberg mass, 
parametrically $\sim g^2 f/16\pi^2 m$). In the sBootstrap bloom, 
the $b$-direction acquires the \emph{largest} mass 
($m_b \gg m_s$ and $m_b > m_c$). The bloom is far more dramatic than a 
perturbative CW lifting---it is a nonperturbative, strong-dynamics 
effect in which confinement generates a mass of order $\Lambda_{\text{QCD}}^2/m_q$, 
not $g^2/(16\pi^2)$. Numerically, the one-loop CW mass for the
pseudo-modulus at the Seiberg vacuum is $m_{\text{CW}} \approx 13\MeV$,
more than two orders of magnitude below the physical $m_b = 4{,}180\MeV$.
This inversion is a signature of the dynamical (SQCD) realization
versus the perturbative (superpotential) one.

\paragraph{A candidate mechanism: fractional instantons.}
On $\mathbb{R}^3 \times S^1$, SU($N_c$) instantons fractionalize into
$N_c$ monopole-instantons \cite{Unsal2008}. Each monopole-instanton
carries a single fermionic zero mode per flavor; when the zero mode is
lifted by the quark mass $m_i$, the amplitude acquires a factor
$\sqrt{m_i}$ (not $m_i$). Since the $N_c = 3$ monopole contributions
enter \emph{additively} in the effective superpotential,
\begin{equation*}
  W_{\text{mon}} \sim \zeta\,\bigl(\sqrt{m_1} + \sqrt{m_2} + \sqrt{m_3}\bigr)\,,
\end{equation*}
where $\zeta \propto e^{-S_0/3}$ is the monopole fugacity. This is
\emph{exactly} the sum-of-square-roots structure of the vacuum charge
$v_0 = (\sum \sqrt{m_k})/3$.
The $v_0$-doubling then maps to the counting of active
monopole-instanton contributions: the seed (two active flavors) involves
two monopole amplitudes, while the bloom (three active flavors, one with
a sign flip) involves three, doubling $v_0$.

This mechanism is non-holomorphic in the sense required: the
K\"ahler potential receives monopole--anti-monopole (bion) cross terms
of the form $|\!\sum s_k\sqrt{m_k}|^2$, where $m_k$ are the
current quark masses. Expanding for $N_c = 3$:
\begin{equation*}
  K_{\text{bion}} = \frac{\zeta^2}{\Lambda^2}\bigl(
  m_1 + m_2 + m_3
  + 2s_1 s_2\sqrt{m_1 m_2}
  + 2s_1 s_3\sqrt{m_1 m_3}
  + 2s_2 s_3\sqrt{m_2 m_3}\bigr)\,,
\end{equation*}
where the diagonal terms $s_k^2 m_k = m_k$ are self-bions and the
cross terms $\sqrt{m_i m_j}$ are mixed bion amplitudes---their
geometric mean structure is what produces $\sum\sqrt{m_k}$.
Note that $m_k$ here are the quark masses (superpotential couplings),
\emph{not} the meson VEVs $M^k{}_k \propto 1/m_k$ of the Seiberg
seesaw; the bion amplitude involves quark zero modes, not meson fields.
Organizing the bion amplitudes into
seed-sector ($S_{\text{seed}} = \sqrt{m_s} + \sqrt{m_c}$) and
bloom-sector ($S_{\text{bloom}} = -\!\sqrt{m_s} + \sqrt{m_c} + \sqrt{m_b}$)
channels, the K\"ahler correction takes the form
\begin{equation*}
  \delta K \sim \lambda_1\,|S_{\text{seed}}|^2
  + \lambda_2\,|S_{\text{bloom}} - 2\,S_{\text{seed}}|^2\,.
\end{equation*}
The second term is a perfect square that vanishes
\emph{at the minimum} when $S_{\text{bloom}} = 2\,S_{\text{seed}}$,
i.e., $\sqrt{m_b} = 3\sqrt{m_s} + \sqrt{m_c}$---the
$v_0$-doubling condition. The relation between bion couplings
($\lambda_3 = -4\lambda_2$ for cross-channel bions) is what
enforces the doubling. Numerically,
$S_{\text{bloom}} - 2\,S_{\text{seed}} = 0.023\MeV^{1/2}$
(effectively zero: $m_b = 4{,}177$ vs.\ PDG $4{,}180$).
The Hessian of $\lambda_2\,|S_{\text{bloom}} - 2\,S_{\text{seed}}|^2$
has rank~1: its single nonzero eigenvector is
$\partial f/\partial m_i$ (the gradient of
$f = \sqrt{m_b} - 3\sqrt{m_s} - \sqrt{m_c}$), while the two flat
directions correspond to independent variation of $m_s$ and $m_c$
at fixed $f$.  Thus the bion K\"ahler constrains exactly \emph{one}
mass combination; the remaining two masses are fixed by the Koide
seed in the superpotential. Including propagated uncertainties
($\delta m_s = 0.8\MeV$, $\delta m_c = 20\MeV$), the prediction
is $m_b = 4{,}177 \pm 40\MeV$, consistent with the PDG value
$4{,}180 \pm 30\MeV$ at $0.06\sigma$.

The principal caveat is that this analysis
applies on $\mathbb{R}^3 \times S^1$; connecting to $\mathbb{R}^4$
physics requires \"Unsal's continuity conjecture~\cite{Unsal2012},
which posits that center-symmetric confinement at small $S^1$ is
continuously connected to the large-$S^1$ confining vacuum.
The conjecture is supported by lattice evidence for $SU(N)$ pure
Yang--Mills and by the absence of phase transitions in the
center-stabilized theory, but remains unproven for theories with
matter.

The identification $\langle F_{\Phi_0}\rangle = f \sim \fpi^2$ is the
sBootstrap's central claim: the SUSY-breaking $F$-term is the chiral
condensate parameter.

In SQCD, this is realized dynamically rather than through a 
superpotential: the chiral condensate $\qqbar \neq 0$ acts as the 
nonzero $F$-term, the impossibility of restoring chiral symmetry 
in the QCD vacuum is the dynamical analogue of $f \neq 0$, and the 
pion---the would-be Goldstone boson that acquires mass from explicit 
chiral breaking---plays the role of $\Phi_0$ whose mass is lifted 
by loop corrections (here, by quark masses).

\subsection{The Goldstino direction and the Volkov--Akulov realization}
\label{subsec:va}

When SUSY is spontaneously broken, a massless fermion appears: the 
Goldstino, the fermionic analogue of the Goldstone boson. Its decay 
constant $f_{\text{VA}}$ sets the SUSY-breaking scale through the 
Volkov--Akulov Lagrangian \cite{VolkovAkulov1973}:
\begin{equation}
\mathcal{L}_{\text{VA}} = -f_{\text{VA}}^2 \det\!\left(
\delta_\mu^\nu + \frac{i}{f_{\text{VA}}^2}\,\bar{\lambda}\bar{\sigma}^\nu
\partial_\mu\lambda\right),
\end{equation}
where $\lambda$ is the Goldstino field. At low energies, the Goldstino 
couples derivatively with strength $\sim 1/f_{\text{VA}}^2$, exactly 
as the pion couples with strength $\sim 1/f_\pi^2$ in the chiral 
Lagrangian.

The sBootstrap identification is:
\begin{equation}
f_{\text{VA}} = f_\pi\,, \qquad 
\text{Goldstino} = \text{muon (partner of }\pi\text{)}\,.
\end{equation}
More precisely, the muon is a \emph{pseudo-Goldstino}: in the 
limit $m_q \to 0$ where the pion becomes exactly massless (a true 
Goldstone boson), its fermionic partner also becomes massless (a 
true Goldstino). The finite quark masses explicitly break both 
chiral symmetry (giving $m_\pi \neq 0$) and the sBootstrap SUSY 
(giving $m_\mu \neq m_\pi$), with the splitting controlled by 
$f_\pi$:
\begin{equation}
m_\pi^2 - m_\mu^2 = f_\pi^2\,.
\end{equation}

This is the Volkov--Akulov relation for the pseudo-Goldstino: the 
mass splitting between the scalar (pion) and its fermionic partner 
(muon) is set by the SUSY-breaking scale squared.

The VA mechanism is not a separate breaking from the F-term: it is 
the \emph{low-energy, composite-level description} of the same 
breaking. At the microscopic level, $\langle F\rangle \neq 0$ 
in the meson superfield; at the composite level, this manifests 
as the Goldstino coupling with $f_{\text{VA}} = f_\pi$. These 
two quantities---the Goldstino decay constant and the pion decay 
constant---are \emph{identified} in the sBootstrap; this is the 
central claim, not a derivation.

The identification should not be confused with the standard relation
between the $F$-term VEV and the chiral condensate. In standard
SQCD,
%
\begin{equation*}
\langle F\rangle \sim \frac{|\qqbar|}{\fpi} \sim
\frac{(250\MeV)^3}{92\MeV} \approx 1.7\times 10^5\MeV^2\,,
\end{equation*}
%
while $\fpi^2 \approx 8{,}500\MeV^2$---a factor of ${\sim}20$
discrepancy. The GOR relation \cite{GOR1968},
$m_\pi^2 \fpi^2 = -(m_u + m_d)\qqbar$, connects these through the
small quark masses: the ratio $|\qqbar|/\fpi \sim 20$ is exactly
$(m_u + m_d)^{-1} m_\pi^2 \fpi \approx 20$, so the discrepancy is
the chiral suppression factor, not a tuning. The sBootstrap
identification $f_{\text{VA}} = f_\pi$ is therefore a
\emph{low-energy effective} statement: the breaking scale appearing
in the meson--lepton mass relation is $f_\pi$ (the Goldstone decay
constant), not the microscopic condensate $|\qqbar|^{1/3}$.

The VA realization gives more than a mass formula---it organizes the 
entire breaking narrative. In the chiral limit ($m_q \to 0$), both 
pion and muon are exactly massless: the pion as the Goldstone boson 
of broken chiral symmetry, the muon as the Goldstino of broken 
sBootstrap SUSY. This is the \emph{exact seed}: the $B$-meson 
direction sits at the flat direction ($\zch = 0$), and the 
$\pi$--$\mu$ pair is massless. Turning on quark masses ($m_q \neq 0$) 
does three things simultaneously:
\begin{enumerate}
\item Gives the pion its GOR mass (lifting the flat direction for $\pi$).
\item Gives the muon its mass (the pseudo-Goldstino acquires mass 
when the explicit breaking parameter $m_q$ is nonzero).
\item Generates the $B$-meson mass (the flat direction for $\Phi_0$ 
is lifted by the Coleman--Weinberg potential, i.e., by the same 
quark masses that feed the GOR formula).
\end{enumerate}
The splitting between (1) and (2) is $f_\pi^2$---the VA relation. 
The seed-to-bloom transition of \S\ref{sec:two_tuples} is the same 
physics seen from the charge-space perspective: $m_q \neq 0$ lifts 
the zero, generates the $b$/$B$ mass, flips the sign of $\zch_s$ 
and $\zch_\pi$, and produces the full triple. The VA equation 
\eqref{eq:pimu} is not one fact among many; it is the generating 
relation from which the bloom follows.

A note on quantum numbers: the Goldstino of spontaneously broken 
spacetime SUSY is a neutral fermion (the fermionic component of the 
superfield whose $F$-term breaks SUSY). The muon carries electric 
charge and lepton number. The sBootstrap SUSY, however, is not a 
spacetime symmetry but an emergent hadronic symmetry---a 
Miyazawa-type superalgebra \cite{Miyazawa} acting on composite 
states. In this context, the ``Goldstino'' need not carry the 
quantum numbers of a spacetime supercharge; it carries the quantum 
numbers appropriate to the broken hadronic generator. The muon's 
charge is inherited from the pion multiplet it partners.

\subsection{D-term breaking: electromagnetic splittings}
The unbroken $U(1)_{\text{em}}$ gauge symmetry allows a 
Fayet--Iliopoulos (FI) term \cite{FayetIliopoulos}:
\begin{equation}
\mathcal{L}_{\text{FI}} = \xi\,D_{\text{em}}\,,
\end{equation}
which shifts scalar masses by $\delta m^2 = Q_{\text{em}}\,\xi$, 
where $Q_{\text{em}}$ is the electric charge. This is the 
supersymmetric version of the electromagnetic mass difference.

In QCD without SUSY, the electromagnetic mass splitting of the 
pion is given by the Das--Guralnik--Mathur--Low--Young (DGMLY) 
formula \cite{DGMLY1967}:
\begin{equation}
m_{\pi^\pm}^2 - m_{\pi^0}^2 = \frac{3\alpha_{\text{em}}}{4\pi}\,
m_\rho^2
\approx (35.5\MeV)^2\,,
\end{equation}
which in the SUSY framework becomes a D-term contribution:
\begin{equation}
\xi \sim \frac{3\alpha_{\text{em}}}{4\pi}\,m_\rho^2 
\sim (35\MeV)^2\,.
\end{equation}

In charge space, the D-term shifts $\zch$ by an amount proportional 
to $Q_{\text{em}}$:
\begin{equation}
\zch \to \zch + \delta\zch\,, \qquad 
\delta\zch \approx \frac{Q_{\text{em}}\,\xi}{2\zch}\,.
\end{equation}
For the pion ($Q_{\text{em}} = \pm 1$, $\zch \approx 11.8$): 
$\delta\zch \approx 0.2\MeV^{1/2}$, matching the observed 
$\zch_{\pi^\pm} - \zch_{\pi^0} = 11.82 - 11.62 = 0.20\MeV^{1/2}$.

The D-term effect is subleading to the F-term: 
$\xi \sim (35\MeV)^2$ versus $f_\pi^2 \sim (92\MeV)^2$. It 
operates \emph{within} each supermultiplet entry, splitting charged 
from neutral mesons, while the F-term splits bosons from fermions 
across the supermultiplet.

Witten's inequality \cite{Witten1983} guarantees 
$m_{\pi^\pm} > m_{\pi^0}$ in QCD with electromagnetism---the 
charged pion is always heavier. In the SUSY framework this 
becomes the statement that the D-term contribution is positive 
for $Q_{\text{em}} \neq 0$, consistent with the FI sign.

\subsection{Mapping the breaking to the tuples}
The tuple structure maps onto a layered breaking with two independent 
mechanisms (F-term and D-term) plus a separate heavy-sector origin. 
The Volkov--Akulov realization of \S\ref{subsec:va} is not a third 
mechanism: it is the low-energy, composite-level description of the 
same F-term breaking. The VA relation \eqref{eq:pimu} \emph{sets the 
scale} ($f_{\text{VA}} = f_\pi$); the O'Raifeartaigh/SQCD dynamics 
of \S\ref{subsec:fterm} \emph{provides the mechanism} 
($\langle F\rangle \neq 0$).

\paragraph{Layer 1: F-term/VA (chiral condensate, scale $\sim f_\pi$).}
The chiral condensate $\qqbar \neq 0$ breaks both chiral symmetry 
and the sBootstrap SUSY. This is the primary breaking, affecting 
the light sector. The seed triples $(0, s, c)$ and $(0, \pi, D_s)$ 
are the pre-breaking configurations; the massless member is the flat 
direction whose $F$-term is forced nonzero by confinement. After 
breaking:
\begin{itemize}
\item The zero member is lifted and acquires negative charge 
($\zch < 0$), becoming the sign-flipped member of the full triple 
($-s$, $-\pi$).
\item The full triples $(-s,c,b)$ and $(-\pi, D_s, B)$ emerge, 
incorporating the $b$-quark and $B$-meson sectors.
\item The lepton triple $(e, \mu, \tau)$ is the fermionic shadow: 
the leptons acquire masses through the same condensate, with the 
muon as pseudo-Goldstino.
\end{itemize}

\paragraph{Layer 2: D-term (electromagnetism, scale $\sim \alpha_{\text{em}} m_\rho$).}
The $U(1)_{\text{em}}$ FI term generates electromagnetic mass 
splittings within each multiplet. This is a perturbative correction 
on top of the F-term breaking, affecting only charged particles. 
It does not create new Koide triples but shifts existing ones by 
$O(\alpha_{\text{em}})$.

\paragraph{Layer 3: The heavy sector (scale $\sim v_0^{(cbt)}$).}
The $(c,b,t)$ triple at $v_0^2 \approx 29.6\GeV$ operates far 
above the chiral scale. It has no seed (no massless member), no 
meson shadow (top does not hadronize), and its breaking may be 
connected to a different sector of the theory---possibly the 
electroweak Higgs mechanism.

This layered structure involves three distinct scales, each with 
different physical meaning:
\begin{align}
\text{F-term (boson--fermion splitting):} &\quad 
f_\pi \approx 92\MeV\,, \quad f_\pi^2 \approx 8{,}500\MeV^2\,; \\
\text{D-term (EM fine structure):} &\quad 
\xi^{1/2} \approx 35\MeV\,, \quad \xi \approx 1{,}200\MeV^2\,; \\
\text{Heavy sector (Koide scale):} &\quad 
v_0^{(cbt)} \approx 172\MeV^{1/2}\,, \quad (v_0^{(cbt)})^2 \approx 29{,}600\MeV\,.
\end{align}
Comparing the squared scales: $f_\pi^2 \gg \xi$ (F-term dominates 
over D-term by a factor $\sim 7$, consistent with the D-term being 
a perturbative $O(\alpha_{\text{em}})$ correction). The heavy-sector 
scale $(v_0^{(cbt)})^2$ carries dimensions of $\MeV$ (not $\MeV^2$) 
because $v_0$ is a charge ($\MeV^{1/2}$), not an energy; direct 
comparison with $f_\pi$ requires specifying which physical quantity 
is being compared.

\subsection{The O'Raifeartaigh structure in SQCD}

In Seiberg's SQCD with $N_f = N_c = 3$, the quantum constraint 
\eqref{eq:seiberg} provides a natural setting for the 
O'Raifeartaigh mechanism. The constraint $\det M - B\bar{B} = \Lambda^6$ 
means the moduli space has a specific geometry: not all field directions 
are flat. When quark masses are added via a tree-level superpotential 
$W_{\text{tree}} = m_i M^i{}_i$ (sum over flavors), the $F$-terms become:
\begin{equation}
F_{M^i{}_j} = m_i\,\delta^i_j + \Lambda^6\,(\text{constraint terms})\,.
\end{equation}

For the seed triple $(0, s, c)$, the ``$0$'' member corresponds to a 
direction in the meson matrix where the tree-level mass vanishes---the 
entry involving the lightest quark combination. In the exact chiral 
limit ($m_u = m_d = 0$), this entry is truly massless, and the 
constraint \eqref{eq:seiberg} forces $\langle F\rangle \neq 0$ for 
that entry. This is the dynamical O'Raifeartaigh mechanism:
confinement (encoded in $\Lambda^6$) plays the role of the parameter
$f$ in the superpotential, and the impossibility of setting all
$F$-terms to zero is a consequence of the strong dynamics rather
than a choice of superpotential.

\paragraph{The pseudo-modulus and the Koide seed.}
In the ISS magnetic dual, the tree-level superpotential contains a
Yukawa coupling $h\,\mathrm{Tr}(q\,M\,\tilde q)$ between the magnetic
quarks $q$, $\tilde q$ and the meson matrix $M$. The O'Raifeartaigh
parameters map as $g \to h$, $m \to h\mu$, $v \to \langle X\rangle$
(the pseudo-modulus VEV), so the seed Koide condition
$gv/m = \sqrt{3}$ becomes $\langle X\rangle/\mu = \sqrt{3}$---the
magnetic Yukawa~$h$ cancels. The Koide seed fixes the
pseudo-modulus VEV, not the coupling~$h$.
At one-loop Coleman--Weinberg level $\langle X\rangle = 0$; a
non-canonical K\"ahler correction provides the mechanism.

\paragraph{K\"ahler stabilization of the pseudo-modulus.}
In the minimal O'Raifeartaigh model ($W = fX + m\phi\tilde\phi
+ gX\phi^2$) with canonical K\"ahler potential,
the CW potential is monotonically increasing and selects
$\langle X\rangle = 0$.  Polynomial superpotential deformations
(including cubic, quartic, and dimension-five operators)
either vanish at the vacuum, reinforce $v = 0$, or produce
minima at $gv/m \ll \sqrt{3}$.

A negative quartic K\"ahler correction resolves this:
\begin{equation*}
  K = |X|^2 + \frac{c}{\Lambda_K^2}\,|X|^4\,,
  \qquad c < 0\,.
\end{equation*}
The tree-level potential becomes
$V_{\text{tree}} = f^2/(1 + 4c\,v^2/\Lambda_K^2)$, which has
a pole at $v_{\text{pole}} = \Lambda_K/(2\sqrt{|c|})$.
Near the pole the CW contribution is negligible
($V_{\text{CW}} \sim 10^{-4}\,m^4 \ll V_{\text{tree}}$),
and the effective potential is minimized at $v_{\text{pole}}$.
Demanding $v_{\text{pole}} = \sqrt{3}\,m/g$ gives the exact condition
\begin{equation*}
  c = -\,\frac{\Lambda_K^2}{12\,m^2}\,.
\end{equation*}
With $\Lambda_K = m$ this is $c = -1/12$, an algebraic result:
$\Lambda_K/(2\sqrt{1/12}) = \sqrt{12}/2 = \sqrt{3}$.
No numerical tuning is required.
At this VEV the fermion spectrum is
$(0,\, (2{-}\sqrt{3})\,m,\, (2{+}\sqrt{3})\,m)$
with $Q = 4/6 = 2/3$ exactly
(since $(\sqrt{2{-}\sqrt{3}} + \sqrt{2{+}\sqrt{3}})^2 = 4 + 2 = 6$),
confirming the Koide seed. The K\"ahler metric
$K_{X\bar X} = 1 - v^2/(3m^2)$ vanishes at $v = \sqrt{3}\,m$,
signaling the boundary of the effective field theory.

The ISS mechanism (Intriligator--Seiberg--Shih \cite{ISS2006}) 
establishes metastable SUSY-breaking vacua for $N_c < N_f < 3N_c/2$, 
where the rank condition in the magnetic dual prevents all $F$-terms 
from vanishing. The sBootstrap case $N_f = N_c = 3$ is strictly 
\emph{below} the ISS window---at $N_f = N_c$ the low-energy 
description is qualitatively different, governed by the 
quantum-modified constraint \eqref{eq:seiberg} on the moduli space 
rather than by a free magnetic dual. However, physical QCD has
$N_f = 3$ light flavors and $N_f = 4$--$5$ if charm and bottom are
included at their respective thresholds. With charm included,
$N_f = 4$ falls within the ISS window ($3 < 4 < 4.5$); with both
charm and bottom, $N_f = 5$ falls in the conformal window
($4.5 < 5 < 9$), where the Seiberg seesaw constraint does not apply
directly. The sBootstrap's multi-scale structure---with
the seed triples involving $s$, $c$, and $b$ quarks at different
scales---requires treating the effective $N_f$ as
scale-dependent: $N_f = 3$ below $m_c$ (confining, Seiberg
constraint), $N_f = 4$ above $m_c$ (ISS metastable SUSY breaking),
and $N_f = 5$ above $m_b$ (conformal window).
Note that $N_f = 4$ is the \emph{unique} integer satisfying
$3 < N_f < 4.5$ for $N_c = 3$, providing a dynamical rationale for
exactly four light flavors $(u,d,s,c)$ in the SUSY-breaking sector.

\paragraph{ISS vacuum lifetime.}
The metastable ISS vacuum decays to the SUSY-preserving ground state
by tunneling with bounce action
$S_{\text{bounce}} \sim (8\pi^2/3)\,N_c\,(\Lambda_{\text{mag}}/h\mu)^4$,
where $\mu$ is the lightest quark mass controlling the tunneling.
For $\mu = m_s = 93\MeV$, $h \sim O(1)$, and
$\Lambda_{\text{mag}} \sim 300\MeV$:
$S_{\text{bounce}} \sim 10^2$--$10^3$ (depending on~$h$),
giving a lifetime exceeding the age of the universe by hundreds of
orders of magnitude for any $h$ of order unity.
The metastable vacuum is cosmologically acceptable.

\paragraph{The Seiberg seesaw and dual Koide.}
At the Seiberg vacuum with $B = \bar B = 0$, the diagonal meson VEVs
scale as $\langle M^j{}_j\rangle \propto 1/m_j$: heavy quarks produce
light mesons and vice versa. This seesaw prompts a dual question: does
the inverse-mass spectrum $\{1/m_i\}$ satisfy a Koide condition?
Evaluating $Q = \sum(1/m_i)/(\sum\sqrt{1/m_i})^2$ for the pure
down-type quarks $(d,s,b)$:
\begin{equation*}
  Q(1/m_d,\, 1/m_s,\, 1/m_b) = 0.6652\,,
\end{equation*}
a 0.22\% deviation from $2/3$---better than any direct quark Koide
triple. No other quark triple comes close under
the Seiberg seesaw (the next best, $(d,s,c)$, deviates by 4.2\%).
This ``dual Koide'' on $(d,s,b)$ (first reported here) connects
the UV mass hierarchy to the IR meson VEVs: the Seiberg constraint
maps the down-type quark masses into a near-Koide spectrum of
meson condensates. However, the
result is sensitive to $m_d$: within $\pm 1\sigma$ ($m_d = 4.67
\pm 0.48\MeV$), $Q$ ranges from $0.655$ to $0.676$, and the deviation
from $2/3$ varies from $0.2\%$ to $1.7\%$. The dual Koide is an
observation at the edge of the current $m_d$ uncertainty.

A simple algebraic identity connects the seed to the seesaw:
for any two masses, the Koide ratio
$Q(a,b) = (a+b)/(\sqrt{a}+\sqrt{b})^2$ satisfies
$Q(a,b) = Q(1/a,\,1/b)$ identically.
The seed $(0, m_s, m_c)$ sits at this self-dual point:
its Koide ratio is the same whether evaluated on quark masses or
on Seiberg-dual meson VEVs $\propto 1/m_i$.  Thus the seed Koide
in the electric (UV) theory is \emph{automatically} a seed Koide
in the magnetic (IR) theory---the Seiberg duality preserves the
seed condition without fine-tuning.
For the full bloomed triple $(-s,c,b)$, self-duality breaks:
$Q(-s,c,b) \neq Q(1/m_s,1/m_c,1/m_b)$ in general.
The dual Koide $Q(1/m_d,1/m_s,1/m_b) = 0.665$ is therefore a
\emph{nontrivial} constraint on the down-type sector, not a
tautological restatement of the seed.
Among all $\binom{5}{3} = 10$ quark triples, $(d,s,b)$ is the
unique triple with inverse-mass Koide ratio within 1\% of $2/3$;
the next closest is $(d,s,c)$ at 4.2\% deviation---a 19:1 margin.

\subsection{How the three supermultiplets join}

The seed-to-bloom transition generates three separate Koide triples. 
But the full sBootstrap requires these triples to be \emph{embedded 
in a single meson matrix}. The $3\times 3$ light-flavor matrix is:

\begin{equation}
M^{(3)} = \begin{pmatrix}
u\bar{u} & u\bar{d} & u\bar{s} \\
d\bar{u} & d\bar{d} & d\bar{s} \\
s\bar{u} & s\bar{d} & s\bar{s}
\end{pmatrix}
\;\longleftrightarrow\;
\begin{pmatrix}
\pi^0/\eta & \pi^+ & K^+ \\
\pi^- & \pi^0/\eta & K^0 \\
K^- & \bar{K}^0 & \eta'
\end{pmatrix}.
\end{equation}

When charm and bottom are included, the matrix extends to $5\times 5$. 
We display it schematically, with the members of the meson Koide 
triple $(-\pi, D_s, B)$ marked by boxes:

\begin{equation}
M^{(5)} = \begin{pmatrix}
\pi^0/\eta & \boxed{\pi^+} & K^+ & D^0 & \boxed{B^+} \\
\pi^- & \pi^0/\eta & K^0 & D^- & B^0 \\
K^- & \bar{K}^0 & \eta' & \boxed{D_s^-} & B_s \\
\bar{D}^0 & D^+ & D_s^+ & \eta_c & B_c \\
B^- & \bar{B}^0 & \bar{B}_s & \bar{B}_c & \eta_b
\end{pmatrix},
\end{equation}
where rows are indexed by $(u, d, s, c, b)$ and columns by 
$(\bar{u}, \bar{d}, \bar{s}, \bar{c}, \bar{b})$.

The three members of the meson Koide triple occupy different 
off-diagonal blocks of this matrix: $\pi^+ = u\bar{d}$ is in the 
light $(u,d)$ block, $D_s^- = s\bar{c}$ bridges the light-charm 
sector, and $B^+ = u\bar{b}$ bridges the light-bottom sector. They 
do not form a ``slice'' in any simple linear-algebraic sense---they 
are scattered entries connected by the Koide condition on their 
\emph{masses}, not by their positions in the matrix. The meson seed 
$(0, \pi, D_s)$ drops the $B^+$ and replaces it with zero (the 
entry that would involve the top quark, absent from the matrix).

This scattered structure is the central puzzle: the Koide triples 
are \emph{mass relations} among matrix entries, not submatrices. 
The question is whether the Seiberg constraint 
$\det M - B\bar{B} = \Lambda^6$ (appropriately generalized to 
$5 \times 5$), combined with the SUSY-breaking dynamics, produces 
correlations among these scattered entries that enforce the 
energy-balance condition $\langle z_k^2\rangle = z_0^2$.

\subsection{The energy balance is not an RG attractor}
\label{subsec:fixedpoint}
One might hope that the energy-balance condition
$\langle z_k^2\rangle = z_0^2$ is a \emph{dynamical attractor}:
that three masses evolving under SUSY-breaking dynamics flow toward
the Koide surface $Q = 2/3$. However, the one-loop soft-mass RG
equations in $\mathcal{N}=1$ SQCD with universal Casimirs and
negligible Yukawas give $dm_i/dt = \gamma(M_g^2 - m_i)$, which
drives all masses to a common value $m_i \to M_g^2$, yielding
$Q \to 1/3$---the \emph{opposite} extreme. The Koide preservation
condition $\sum_i F_i(3 - 2S/\sqrt{m_i}) = 0$ on the $Q = 2/3$
surface is algebraically incompatible with any uniform attractor
flow (by the Cauchy--Schwarz inequality, the required identity
$S\sum 1/\sqrt{m_i} = 9/2$ contradicts $S\sum 1/\sqrt{m_i} \geq 9$).

The energy-balance condition therefore does \emph{not} arise from
simple RG running. It must instead be a \emph{boundary condition}
on the superpotential---as in the O'Raifeartaigh realization of
\S\ref{subsec:fterm}, where $gv = \sqrt{3}\,m$ produces the exact
Koide seed. The condition is set at the scale of SUSY breaking and
then preserved by the QCD-invariant mass ratios, rather than being
an RG attractor.

\subsection{The exact seed and what breaks it}
\label{subsec:exactseed}

The seed configuration $(0, \zch_a, \zch_b)$ requires \emph{two} 
independent conditions:
\begin{enumerate}
\item One member at $\zch = 0$ exactly (the flat direction).
\item The ratio of the massive members equals $2 - \sqrt{3}$:
\begin{equation}
\frac{\sqrt{m_a}}{\sqrt{m_b}} = 2 - \sqrt{3} \approx 0.26795\,.
\end{equation}
\end{enumerate}
These are logically distinct. The chiral limit ($m_u = m_d = 0$, 
$\qqbar \neq 0$) enforces condition (1) for the meson sector: the 
pion is an exact Goldstone boson, $m_\pi = 0$, and in the 
picture of \S\ref{subsec:fterm} the $B$-meson direction sits at 
the flat direction before acquiring mass. But the chiral limit says 
\emph{nothing} about condition (2)---it does not constrain 
$\sqrt{m_{D_s}}/\sqrt{m_B}$ or $\sqrt{m_s}/\sqrt{m_c}$ 
to equal $2 - \sqrt{3}$.

The physical deviations from the exact ratio are small but 
measurable:

\begin{center}
\begin{tabular}{lccc}
\toprule
Seed & Exact ratio & Physical ratio & Deviation \\
\midrule
$(0, \pi, D_s)$ & 0.2680 & 
$\sqrt{m_\pi}/\sqrt{m_{D_s}} = 0.2663$ & $-0.6\%$ \\
$(0, s, c)$ & 0.2680 & 
$\sqrt{m_s}/\sqrt{m_c} = 0.2712$ & $+1.2\%$ \\
$(0 \approx e, \mu, \tau)$ & 0.2680 & 
$\sqrt{m_\mu}/\sqrt{m_\tau} = 0.2439$ & $-9.0\%$ \\
\bottomrule
\end{tabular}
\end{center}

The meson seed is closest to the exact ratio (0.6\% deviation). 
The quark seed deviates by 1.2\%. A natural thought is that QCD 
running might shift the quark ratio: perhaps at some scale $\mu^*$ 
the ratio $\sqrt{m_s(\mu^*)}/\sqrt{m_c(\mu^*)}$ hits $2 - \sqrt{3}$ 
exactly. But this is \emph{not possible within QCD alone}. At leading 
order, all quark masses run with the same anomalous dimension 
$\gamma_m$:
\begin{equation}
m_i(\mu) = m_i(\mu_0)\left(\frac{\alpha_s(\mu)}{\alpha_s(\mu_0)}\right)^{\gamma_m/\beta_0},
\end{equation}
so the ratio $m_s(\mu)/m_c(\mu)$---and therefore 
$\sqrt{m_s}/\sqrt{m_c}$---is an RG invariant of QCD. Higher-order 
corrections break this invariance, but only at the $O(\alpha_s^2)$ 
level, far too small to account for 1.2\%.

The quark seed deviation is therefore \emph{intrinsic}: the 
$\overline{\text{MS}}$ mass ratio $m_s/m_c$ is a fixed number 
(at any QCD scale), and it misses the Koide target by 1.2\%. 
Three possible explanations present themselves:

\paragraph{Electroweak running.} Above the electroweak scale, 
quark masses are controlled by Yukawa couplings $y_i$ to the 
Higgs, which \emph{do} run differently for different flavors: 
the $\beta$-functions for $y_s$ and $y_c$ involve different 
electroweak quantum numbers ($SU(2)_L$ doublet structure, 
hypercharge). The ratio $y_s/y_c$ evolves with scale, and 
there may exist a scale $\mu^*$ above the electroweak 
threshold where $\sqrt{y_s(\mu^*)}/\sqrt{y_c(\mu^*)} = 2 - \sqrt{3}$. 
This would identify the ``seed scale'' as an electroweak-sector 
phenomenon, not a QCD one---consistent with the sBootstrap's 
identification of the breaking with $SU(2)_L \times U(1)_Y$ 
rather than $SU(3)_c$.

\paragraph{Charge running.} Even if mass ratios are QCD-invariant,
the Koide decomposition $\zch_k = z_0 + z_k$ may have its own
flow. The Koide scale $v_0 = z_0^2$ sets the common scale,
and the fluctuations $z_k$ encode the deviations. As shown in
\S\ref{subsec:fixedpoint}, simple RG running drives $Q \to 1/3$
rather than $2/3$, so the energy-balance condition cannot be an
attractor in the usual sense. However, the O'Raifeartaigh
realization (\S\ref{subsec:fterm}) shows that the condition is a
boundary condition on the superpotential ($gv = \sqrt{3}\,m$).
The 1.2\% deviation of the quark seed ratio may then reflect
the distance between the physical quark masses and the exact
superpotential constraint.

\paragraph{Meson vs.\ quark deviations.} The meson deviation 
(0.6\%) is smaller than the quark deviation (1.2\%), and for 
a different reason: meson masses are \emph{not} simple functions 
of quark mass ratios. The pion mass is given by the GOR 
relation $m_\pi^2 f_\pi^2 = -(m_u + m_d)\qqbar$; the $D_s$ 
mass involves heavy-quark effective theory; and both receive 
chiral perturbation theory corrections. These nonperturbative 
QCD effects are the dominant source of the meson seed deviation, 
and are in principle calculable at NLO in ChPT.

\paragraph{A remarkable feature of the VA relation.} The 
foundational equation $m_\pi^2 - m_\mu^2 = f_\pi^2$ contains 
no QCD coupling constant: no $\alpha_s$, no $N_c$, no 
$\Lambda_{\text{QCD}}$. It is a pure mass relation among three 
experimentally measured quantities---two masses and a decay 
constant---holding to 2.2\%. In the sBootstrap, this is natural: 
the relation is a SUSY mass splitting, set by the breaking 
scale $f_\pi$, and SUSY mass relations do not involve the gauge 
coupling of the confining theory. Standard QCD \emph{explains} 
$f_\pi$ in terms of $\qqbar$ and $\alpha_s$; the sBootstrap 
\emph{uses} $f_\pi$ as the fundamental parameter, with QCD 
providing the dynamical origin but not appearing in the mass 
formula. The absence of QCD factors in \eqref{eq:pimu} is a 
hint that the relevant symmetry is not $SU(3)_c$ but the 
hadronic superalgebra that connects pions to leptons.

The lepton ``seed'' has the worst ratio match ($-9\%$), which 
is consistent with the observation that the lepton triple is 
\emph{not} a true seed: $m_e \neq 0$, and the electron was never 
at the flat direction. The leptons are an ``almost-seed'' whose 
excellent $Q = 0.66666$ arises not from the seed structure but 
from the smallness of $m_e$, which makes the two conditions 
approximately degenerate. For the leptons, condition (1) is 
approximate rather than exact, and condition (2) is correspondingly 
relaxed.

\subsection{What the $(c,b,t)$ triple tells us}

The $(c,b,t)$ triple ($Q = 0.6695$, $v_0^2 = 29.6\GeV$) has 
neither a seed nor a meson shadow. It operates at a scale far 
above $\Lambda_{\text{QCD}}$, suggesting its origin is different:

\begin{enumerate}
\item It may be governed by the electroweak Higgs mechanism rather
than the chiral condensate. In a two-Higgs-doublet model with
$\tan\beta = 1$ (as predicted by the mixed-triple Koide condition,
\S\ref{subsec:superpotential}), $y_t \approx \sqrt{2}$;
the Koide condition on $(c,b,t)$ would then constrain the Yukawa
coupling matrix.
\item It may represent a Volkov--Akulov realization at the 
electroweak scale: if there is a second SUSY breaking at 
$f' \sim v_{\text{EW}}$, the $(c,b,t)$ triple would be its 
analog of the $(\pi, \mu)$ system, but with $f'$ replacing $f_\pi$.
\item Or it may be a purely kinematic accident---the quark masses 
happen to satisfy the Koide condition without a dynamical reason. 
The 0.42\% precision makes this unlikely but not impossible.
\end{enumerate}

Interpretation (1) is most natural within the sBootstrap: the
$(c,b,t)$ triple involves quarks whose masses are dominated by
Yukawa couplings to the Higgs, not by the chiral condensate.
Since $m_i = y_i\,v_{\text{EW}}/\sqrt{2}$, the Koide condition
$Q = 2/3$ translates into a constraint on the Yukawa couplings
(see \S\ref{subsec:superpotential}):
$\sum y_i/(\sum \sqrt{y_i})^2 = 2/3$.
With $y_t \approx 1$ fixed by the top mass, this determines
$y_b/y_t$ and $y_c/y_t$ through the Koide parametrization.
The chiral sector and the electroweak sector would each
independently produce a Koide-type energy balance, with the charm
and bottom quarks serving as the bridge between the two.

\paragraph{CKM mixing and the Koide triples.}
A striking feature of the quark Koide triples is that they
\emph{obligatorily mix up-type and down-type quarks}: $(-s,c,b)$
contains one up-type ($c$) and two down-type ($s,b$), while
$(c,b,t)$ contains two up-type ($c,t$) and one down-type ($b$). An
exhaustive scan of all $\binom{6}{3} = 20$ quark triples with all
sign choices confirms that \emph{every} triple close to $Q = 2/3$ is
mixed---neither the pure down-type $(d,s,b)$ nor the pure up-type
$(u,c,t)$ comes close ($Q = 0.73$ and $Q = 0.85$ respectively).
This mixing of mass eigenvalues from different Yukawa sectors is
puzzling in the Standard Model, where up-type and down-type masses are
eigenvalues of independent matrices $Y_u$ and $Y_d$. It is natural,
however, in an $SU(5)_f$ flavor-symmetric framework where all five
quarks are treated democratically.

The Koide constraints on the two overlapping triples, together with
the seed condition on $(0,s,c)$, determine $(s,c,b)$ from a single
input ($m_t$). The unconstrained masses are $m_u$ and $m_d$---the
first generation. The Oakes relation $\tan\theta_C \approx
\sqrt{m_d/m_s}$ connects precisely these free parameters to the
Cabibbo angle, suggesting that the CKM matrix encodes the residual
freedom left undetermined by the Koide structure.

If interpretation (2) is also invoked, the sBootstrap has two breaking 
scales:
\begin{equation}
f_{\text{chiral}} = f_\pi \approx 92\MeV\,, \qquad
f_{\text{EW}} \sim v_{\text{EW}} \approx 246\GeV\,,
\end{equation}
and the full spectrum requires both. The light sector (leptons, 
light mesons) is governed by $f_\pi$; the heavy sector (top, and 
by extension charm and bottom through the Koide network) is 
governed by $v_{\text{EW}}$. The charm and bottom quarks, 
appearing in both the $(0,s,c)/(-s,c,b)$ light-sector triples 
and the $(c,b,t)$ heavy-sector triple, are the ``bridges'' 
between the two scales.

\subsection{Toward a Superpotential}
\label{subsec:superpotential}

Consider $\mathcal{N}=1$ SQCD with gauge group $SU(3)_c$ and $N_f=5$ flavors
of chiral superfields $Q^i$, $\bar Q_j$ ($i,j=1\ldots 5$, corresponding to
$u,d,s,c,b$) in the fundamental and antifundamental representations. The
tree-level superpotential is
%
\begin{equation*}
  W_{\text{tree}} = \operatorname{Tr}(\hat m\, M)\,,
  \qquad
  \hat m = \operatorname{diag}(m_u,\, m_d,\, m_s,\, m_c,\, m_b)\,,
\end{equation*}
%
where $M^i{}_j = Q^i\cdot\bar Q_j$ is the $5\times 5$ composite meson
superfield. Since $m_c, m_b \gg \Lambda_{\text{QCD}}$, the charm and
bottom quarks are integrated out at tree level, reducing the effective
theory to $N_f = N_c = 3$ light flavors $(u,d,s)$. In this regime,
Seiberg's exact results give confinement with the composite meson
$M^{(3)}$, baryons $B$, $\bar B$, and the quantum-modified
constraint~\eqref{eq:seiberg}:
%
\begin{equation*}
  \det M^{(3)} - B\bar B = \Lambda_{\text{eff}}^{6}\,,
\end{equation*}
%
where $\Lambda_{\text{eff}}$ is the dynamical scale of the effective
three-flavor theory after threshold matching at $m_c$ and $m_b$.
Implementing this constraint via a Lagrange multiplier $X$ gives the full
effective superpotential
%
\begin{equation*}
  W_{\text{eff}}
  = \operatorname{Tr}(\hat m^{(3)}\, M^{(3)})
  + X\!\left(\det M^{(3)} - B\bar B - \Lambda_{\text{eff}}^{6}\right).
\end{equation*}

The $F$-term equations for the $3\times 3$ light-flavor block read
%
\begin{align*}
  \frac{\partial W_{\text{eff}}}{\partial M^i{}_j}
  &= \hat m_j\,\delta^i{}_j
    + X\,(\det M^{(3)})\,(M^{(3)\,-1})_i{}^j = 0\,,\\[4pt]
  \frac{\partial W_{\text{eff}}}{\partial B}
  &= -X\,\bar B = 0\,,\\[4pt]
  \frac{\partial W_{\text{eff}}}{\partial \bar B}
  &= -X\, B = 0\,.
\end{align*}
%
The baryon equations force either $X=0$ or $B=\bar B=0$. If $X=0$ the meson
equation reduces to $\hat m_j\,\delta^i{}_j=0$, which has no solution for
nonzero quark masses. Hence $B=\bar B=0$ and the constraint becomes
$\det M^{(3)} = \Lambda_{\text{eff}}^6$. The diagonal meson equation then requires
%
\begin{equation*}
  \hat m_j + X\,\frac{\det M^{(3)}}{M^j{}_j} = 0
  \qquad(\text{no sum on } j)\,,
\end{equation*}
%
which determines each diagonal entry of $M^{(3)}$ in terms of the quark masses and
$\Lambda_{\text{eff}}$. For a seed configuration in which one quark mass vanishes---say
$m_u\to 0$---the corresponding $F$-term
$F_{M^1{}_1}=X\,(\det M^{(3)})(M^{(3)\,-1})_1{}^1$ remains nonzero so long as
$\det M^{(3)}\neq 0$, which is guaranteed by the quantum constraint. Supersymmetry
is therefore spontaneously broken, with order parameter
%
\begin{equation*}
  \langle F \rangle \;\sim\; \fpi^2\,.
\end{equation*}

The scalar and fermionic components of $M^i{}_j$ acquire different masses
through this breaking. The scalar component (pseudoscalar meson $q_i\bar q_j$)
picks up an $F$-term contribution absent for the fermionic component (lepton,
via the Miyazawa superalgebra):
%
\begin{align*}
  m^2_{M^i{}_j}\big|_{\text{scalar}}
  &= |m_i + m_j|^2 + f^2\,,\\[3pt]
  m^2_{M^i{}_j}\big|_{\text{fermion}}
  &= |m_i + m_j|^2\,,
\end{align*}
%
where $f^2\approx \fpi^2$. The individual masses $|m_i + m_j|$ are
the tree-level (perturbative) contributions; the physical masses are
dominated by nonperturbative QCD effects that shift both meson and
lepton by the same amount.  The mass \emph{splitting}, however,
is universal across the multiplet:
%
\begin{equation*}
  m^2_{\text{meson}} - m^2_{\text{lepton}} = f^2 \approx \fpi^2\,,
\end{equation*}
%
because $f^2$ is a SUSY-breaking contribution independent of the QCD
dynamics.  This is precisely the Volkov--Akulov
relation~\eqref{eq:pimu}.
For heavier mesons the actual splitting is strongly flavor-dependent
(\S\ref{sec:open}, ``kaon--muon puzzle'').

The top quark, with $m_t\gg\Lambda_{\text{QCD}}$, does not hadronize and
therefore does not participate in the confined meson matrix---the counting
$N_f=5$ is exact. Its mass, and corrections to the heavy-quark sector
($c$, $b$), originate instead from electroweak symmetry breaking.
In $\mathcal{N}=1$ SUSY, holomorphy of the superpotential requires
separate Higgs doublets $H_u$ and $H_d$ for up-type and down-type
quarks, giving
%
\begin{equation*}
  W_{\text{Yukawa}}
  = y^u_{ij}\, H_u\, Q^i\, \bar Q_j
  + y^d_{ij}\, H_d\, Q^i\, \bar Q_j\,,
\end{equation*}
%
with $\langle H_u\rangle = v\sin\beta/\sqrt{2}$ and
$\langle H_d\rangle = v\cos\beta/\sqrt{2}$
($v = 246\GeV$, $\tan\beta = v_u/v_d$).
This is relevant only for quarks whose masses are dominated by the
electroweak condensate rather than the chiral condensate. For the light sector
($u,d,s$) the Yukawa contribution is negligible compared to
$\operatorname{Tr}(\hat m\, M)$. Note that $\hat m$ in the chiral
mass term already \emph{is} the current quark mass generated by the
Higgs mechanism; the two terms are not additive.  The Yukawa
superpotential is written separately to make explicit the dependence
on the Higgs doublets, which are needed for the heavy sector where
$m_i \gg \Lambda_{\text{QCD}}$ and the mass cannot be attributed to
chiral symmetry breaking.

The two-doublet structure has a direct consequence for the Koide
triples. Since the quark Koide triples obligatorily mix up-type
and down-type quarks (\S\ref{subsec:cbt}), the translation from
masses to Yukawa couplings depends on~$\tan\beta$.
For the $(c,b,t)$ triple, with $c$ and $t$ receiving mass
from $H_u$ and $b$ from~$H_d$:
%
\begin{equation*}
  Q(m_c, m_b, m_t)
  = Q(y_c\sin\beta,\; y_b\cos\beta,\; y_t\sin\beta)\,,
\end{equation*}
%
which equals $Q(y_c, y_b, y_t)$ only when $\sin\beta = \cos\beta$,
i.e., $\tan\beta = 1$. The same conclusion holds for the $(-s,c,b)$
triple, where $s$ and $b$ come from $H_d$ and $c$ from $H_u$.
Thus $\tan\beta = 1$ is the unique value at which the Koide condition
on \emph{masses} is equivalent to a Koide condition on
\emph{Yukawa couplings} for mixed up/down-type triples.

This prediction is independently motivated by the sBootstrap's
$SU(5)_f$ flavor symmetry, which treats all five quarks
democratically.  A ratio $\tan\beta \neq 1$ would distinguish
up-type from down-type quarks in the vacuum, breaking $SU(5)_f$
at the electroweak scale in a way not present in the chiral sector.
Equal Higgs VEVs, $v_u = v_d = v/\sqrt{2} = 174\GeV$,
preserve the democratic structure.
At $\tan\beta = 1$, the top Yukawa is
$y_t = \sqrt{2}\,m_t/v_u = 2m_t/v \approx 1.40 \approx \sqrt{2}$.

The Koide condition on the $(c,b,t)$ triple constrains the diagonal
Yukawa matrix $Y = \operatorname{diag}(y_c,y_b,y_t)$.
At $\tan\beta = 1$, all Yukawas are
$y_i = 2m_i/v$ regardless of type, and the Koide relation reads
%
\begin{equation*}
  \frac{y_c + y_b + y_t}
       {(\sqrt{y_c}+\sqrt{y_b}+\sqrt{y_t})^2}
  = \frac{2}{3}\,,
\end{equation*}
%
which is equivalent to writing the square roots of the Yukawa
couplings as
%
\begin{equation*}
  \sqrt{y_k} = \sqrt{y_0}\,
  \bigl(1 + \sqrt{2}\,\cos(2\pi k/3 + \delta)\bigr)\,,
  \qquad k = 0,1,2\,,
\end{equation*}
%
where $y_0 = \bigl(\sum\sqrt{y_k}\bigr)^2\!/9$ and $\delta$ is
the Koide phase. For the physical $(c,b,t)$ masses,
$\delta \approx 0.069$ and $Q = 0.6695$ (a 0.42\% deviation from
$2/3$). The traceless part of $\sqrt{Y}$ lies in the Cartan
subalgebra of $SU(3)_f$:
%
\begin{equation*}
  \sqrt{Y} = \sqrt{y_0}\,
  \Bigl(I_3 + \tfrac{\sqrt{6}}{2}\,
  \hat n\cdot\boldsymbol\lambda_{\text{diag}}\Bigr)\,,
\end{equation*}
%
where $\boldsymbol\lambda_{\text{diag}} = (\lambda_3, \lambda_8)$
are the diagonal Gell-Mann matrices and
$\hat n = (\cos\theta,\,\sin\theta)$ is a unit vector
in the Cartan plane.  An exact trigonometric identity gives
$\theta = \pi/6 - \delta$, so the Koide phase is literally
the angle of the Yukawa matrix in flavor space.
The Koide condition fixes the amplitude at $\sqrt{6}/2$; only
$\hat n$ (equivalently $\delta$) is free.
The bloom from seed to physical triple is a rigid rotation in
the $(\lambda_3,\lambda_8)$ plane by $-\Delta\delta$.

The total Lagrangian is
%
\begin{equation*}
  \mathcal{L}
  = \int\!d^4\theta\,K
  + \left[\int\!d^2\theta\,W + \text{h.c.}\right]
  + \mathcal{L}_D
  + \mathcal{L}_{\text{soft}}\,,
\end{equation*}
%
with K\"ahler potential (canonical plus stabilization and bion corrections),
%
\begin{equation*}
  K = \underbrace{\operatorname{Tr}(M^\dagger M) + |X|^2 + |B|^2
    + |\bar B|^2 + |H_u|^2 + |H_d|^2}_{K_{\text{can}}}
  \underbrace{{}-\,\frac{|X|^4}{12\,\mu^2}}_{K_{\text{stab}}}
  + \underbrace{\frac{\zeta^2}{\Lambda^2}\,
    \Bigl|\sum_{k=1}^{N_c} s_k \sqrt{\hat m_k}\,\Bigr|^2}_{K_{\text{bion}}}\,,
\end{equation*}
The stabilization term $K_{\text{stab}} = -|X|^4/(12\mu^2)$ is the
negative quartic K\"ahler correction from \S\ref{subsec:superpotential}
(with $\Lambda_K = \mu$, $c = -1/12$); it pins the pseudo-modulus at
$\langle X\rangle = \sqrt{3}\,\mu$ via the K\"ahler pole.
The monopole fugacity is $\zeta \propto e^{-S_0/(2N_c)}$,
$\hat m_k$ are the current quark masses (the couplings in
$W \ni \operatorname{Tr}(\hat m\,M)$), and
$s_k = \pm 1$ are monopole-instanton signs (encoding the seed-to-bloom
phase). The bion term is non-holomorphic (containing $\sqrt{\hat m}$) and
non-perturbatively suppressed by $e^{-2S_0/3} \sim 10^{-2}$ at the
confinement scale. Its minimum enforces the $v_0$-doubling condition
$\sqrt{m_b} = 3\sqrt{m_s} + \sqrt{m_c}$
(\S\ref{subsec:fterm});
%
superpotential combining mass, constraint, and Yukawa terms,
%
\begin{equation*}
  W = \underbrace{\operatorname{Tr}(\hat m\, M)}_{\text{chiral mass}}
    + \underbrace{X\!\left(\det M^{(3)} - B\bar B
      - \Lambda_{\text{eff}}^{6}\right)}_{\text{Seiberg constraint}}
    + \underbrace{c_3\,\frac{(\det M^{(3)})^3}
      {\Lambda_{\text{eff}}^{18}}}_{\text{3-instanton}}
    + \underbrace{y_c\, H_u\, M^c{}_c
      + y_b\, H_d\, M^b{}_b
      + y_t\, H_u\, Q^t \bar Q_t\phantom{X}}_{\text{Yukawa}}\,,
\end{equation*}
%
where $M^{(3)}$ is the $3\times 3$ light-flavor block ($u,d,s$)
and $\Lambda_{\text{eff}}$ is the dynamical scale of the confined
$N_f = N_c = 3$ theory (including threshold matching at $m_c$
and $m_b$).  The Yukawa terms are written in the full five-flavor
theory above the heavy-quark thresholds, where $M^c{}_c = Q^c\bar Q_c$
and $M^b{}_b = Q^b\bar Q_b$ are elementary bilinears;
below threshold, they reduce to the matching conditions that generate
$\Lambda_{\text{eff}}$ and the Seiberg constraint.
The top Yukawa $y_t\,H_u\,Q^t\bar Q_t$ stands apart because the
top does not confine into a meson.
At $\tan\beta = 1$, the tree-level lightest Higgs mass vanishes
($m_h = m_Z|\cos 2\beta| = 0$); an NMSSM-type coupling
$\lambda\,X\,H_u\!\cdot\!H_d$ lifts the flat direction and gives
$m_h = \lambda v/\sqrt{2} = 125\GeV$ for $\lambda = 0.72$.
The field $X$ (the Seiberg Lagrange multiplier) is a gauge singlet
already present in the superpotential.
The constraint $\tan\beta = 1$ ($v_u = v_d = v/\sqrt{2}$) is imposed
by the Koide condition on mixed triples
(\S\ref{subsec:superpotential}).
The $D$-term Lagrangian includes an optional Fayet--Iliopoulos
parameter $\xi$ for $U(1)_{\text{em}}$:
%
\begin{equation*}
  \mathcal{L}_D
  = -\frac{e^2}{2}\!\left(\sum_i Q_i\,|\phi_i|^2
    + \xi\right)^{\!2}.
\end{equation*}
%
Finally, soft SUSY-breaking terms arise from a spurion
$\langle F_S\rangle = \fpi^2$:
%
\begin{equation*}
  \mathcal{L}_{\text{soft}}
  = -\tilde m^2 \operatorname{Tr}(M^\dagger M)
  - \bigl(B_\mu\operatorname{Tr}(\hat m\, M)
    + A_y\bigl(y_c\, H_u\, M^c{}_c + y_b\, H_d\, M^b{}_b\bigr) + \text{h.c.}\bigr)\,,
  \qquad
  \tilde m^2 \sim \fpi^2\,.
\end{equation*}
%
The soft scalar mass $\tilde m^2 \sim \fpi^2 \approx (92\MeV)^2$
generates the universal meson--lepton splitting
$m^2_{\text{meson}} - m^2_{\text{lepton}} = \fpi^2$.

\paragraph{Scope of the Lagrangian.}
This Lagrangian produces: (i)~the seed Koide condition ($gv/m = \sqrt{3}$
in the O'Raifeartaigh limit); (ii)~the meson--lepton mass splitting
($\fpi^2$); (iii)~the $v_0$-doubling condition via the bion K\"ahler
correction $K_{\text{bion}}$, whose minimum enforces
$\sqrt{m_b} = 3\sqrt{m_s} + \sqrt{m_c}$; (iv)~the Yukawa couplings for
the heavy sector; (v)~the $\cos(9\delta)$ potential that selects
the lepton Koide phase, via the three-instanton operator
$(\det M^{(3)})^3/\Lambda_{\text{eff}}^{18}$. The bloom is
non-holomorphic and lives in the K\"ahler sector: the bion term
$|\sum s_k\sqrt{\hat m_k}|^2$ is the leading monopole-instanton
contribution on $\mathbb{R}^3 \times S^1$
(\S\ref{subsec:fterm}).

\paragraph{Mass spectrum from the Lagrangian.}
The mass of each particle is determined by a specific term:
\begin{center}
\small
\begin{tabular}{lll}
\toprule
Mass & Origin & Equation \\
\midrule
$m_{\text{meson}}^2 - m_{\text{lepton}}^2 = \fpi^2$
  & $\mathcal{L}_{\text{soft}}$: $\tilde m^2 \sim \fpi^2$ & Universal splitting \\
$m_c/m_s = (2+\sqrt{3})^2$
  & $W$: O'Raifeartaigh at $gv/m = \sqrt{3}$ & Seed Koide \\
$\sqrt{m_b} = 3\sqrt{m_s} + \sqrt{m_c}$
  & $K_{\text{bion}}$: minimum of $|S_{\text{bloom}} - 2S_{\text{seed}}|^2$
  & $v_0$-doubling \\
$m_t = y_t v_u/\sqrt{2}$, $\tan\beta = 1$
  & $W_{\text{Yukawa}}$: $y_t H_u Q^t\bar Q_t$ & External \\
$\delta m^2 = Q_{\text{em}}\,\xi$
  & $\mathcal{L}_D$: FI term & EM splitting \\
$\delta_\ell \bmod 2\pi/3 = 2/9$
  & $W$: $(\det M^{(3)})^3/\Lambda^{18}$ & 3-instanton \\
\bottomrule
\end{tabular}
\end{center}
The first four rows account for six quark masses ($m_u$ and $m_d$ free,
$m_s$ sets the scale, $m_c$ from seed, $m_b$ from K\"ahler, $m_t$
external). The meson--lepton splitting is universal across the
multiplet, and the $D$-term provides the subdominant electromagnetic
correction.

The parameters of this construction separate into those fixed by
structure and those that remain free.

\paragraph{Fixed by the Koide seed.}
The O'Raifeartaigh condition $gv/m = \sqrt{3}$ translates, in the
ISS magnetic dual, into $\langle X\rangle = \sqrt{3}\,\mu$ for the
pseudo-modulus.
This fixes the mass ratio $m_c/m_s = (2+\sqrt{3})^2$
within the seed triple, predicting $m_c = 1{,}301\MeV$ from $m_s$ alone
(PDG: $1{,}270\MeV$, 2.4\% deviation).

\paragraph{Fixed by the vacuum-charge doubling.}
The bloom from seed to full triple doubles $v_0$, giving
$\sqrt{m_b} = 3\sqrt{m_s} + \sqrt{m_c}$ and predicting
$m_b = 4{,}177\MeV$ (0.1$\sigma$ from PDG). This constraint
is \emph{non-holomorphic}: it is a linear relation among square
roots of masses, while the superpotential is holomorphic in meson
fields whose VEVs scale as $1/m_i$. The map
$\zch_k = \sqrt{m_k} \sim 1/\!\sqrt{\langle M^k{}_k\rangle}$
is non-holomorphic. The $v_0$-doubling must therefore arise from the
K\"ahler potential, nonperturbative effects, or emergent dynamics---it
is not a tree-level superpotential term.

\paragraph{Fixed by SUSY breaking.}
The supertrace relation $m^2_\pi - m^2_\mu = \fpi^2$ fixes the
splitting scale. At the Seiberg vacuum, $F_X = 0$ (SUSY is
unbroken in the confined sector), but the Yukawa couplings
produce $F_{H_u} = y_c M^c{}_c$ and $F_{H_d} = y_b M^b{}_b$.
By the seesaw $y_j M_j = (m_j/v)(C/m_j) = C/v$, these
F-terms are flavour-universal:
$F_{H_u} = F_{H_d} = C/v_{\text{EW}} = 290\MeV$.
The identification $\langle F\rangle \sim \fpi^2$
links $\Lambda_{\text{QCD}}$ to $m_s$ via the GOR relation.

\paragraph{$D$-term effects on the Koide condition.}
In $\mathcal{N}=1$ SUSY, $D$-terms shift scalar masses but leave
fermion masses untouched. The lepton Koide triples, being conditions
on \emph{fermion} masses, are naturally protected against
$D$-term perturbations. The meson Koide triple
$(-\pi, D_s, B)$ involves scalar masses and \emph{is} in principle
affected by $D$-terms: a FI parameter $\xi$ shifts $m^2 \to m^2 + \xi Q_{\text{em}}^2$.
Since all three mesons carry unit charge, each $m^2$ shifts by the
same $\xi$, and the relative correction to $Q$ is of order
$\xi/(4m_\pi^2) \sim 0.1\%$ for $\xi \approx -80\MeV^2$.
Numerically, the $D$-term shifts the meson $Q$ by $0.03\%$,
well below the $0.10\%$ precision of the triple.

\paragraph{Remaining free parameters.}
After the seed Koide fixes $m_c/m_s$, the $v_0$-doubling fixes
$m_b$, and $\langle X\rangle/\mu = \sqrt{3}$ is imposed, only two to three free
parameters remain:
\begin{itemize}
  \item $m_u$, $m_d$: unconstrained by any Koide triple; connected
    to the Cabibbo angle through the Oakes relation
    $\tan\theta_C \approx \sqrt{m_d/m_s}$ (giving $V_{us} = 0.224$,
    PDG: $0.2243$);
  \item $m_s$ (or equivalently $\Lambda_{\text{QCD}}$): the overall
    scale of the light sector. The lepton Koide scale
    $M_0 \approx \fpi^2/27$ (\S\ref{sec:seeds}) ties the lepton
    scale to $\fpi$ and hence to $\Lambda_{\text{QCD}}$.
\end{itemize}
The top mass $m_t$ is external---it does not participate in the
confined meson matrix---but if $Q(c,b,t) = 2/3$ is imposed as a
constraint rather than treated as a test, $m_t$ is predicted from
$m_s$ alone. The Yukawa couplings
$y_c = \sqrt{2}\,m_c/v_{\text{EW}}$,
$y_b = \sqrt{2}\,m_b/v_{\text{EW}}$ are fixed by the predicted masses.

The $(c,b,t)$ Koide triple provides an \emph{overdetermination test}.
Using the self-consistent predictions
$m_c^{\text{seed}} = 1{,}301\MeV$ and
$m_b^{v_0} = 4{,}233\MeV$ (from the seed and $v_0$-doubling
with input $m_s$ only), together with the external $m_t$:
$Q(c,b,t) = 0.668$, a $0.14\%$ deviation from~$2/3$---closer
than the $0.42\%$ obtained with PDG masses.
The self-consistent predictions satisfy the Koide conditions
better than the measured values.
If instead both overlapping triples are required to
satisfy $Q = 2/3$ exactly, the overlap system determines
$m_b^{\text{pred}} = 4{,}159\MeV$ (0.5\% from PDG) and
$m_c^{\text{pred}} = 1{,}369\MeV$ (7.8\% from PDG). The bion
K\"ahler prediction $m_b = 4{,}177\MeV$ ($0.07\%$, $0.06\sigma$)
is more precise, suggesting that $v_0$-doubling rather than exact
$Q = 2/3$ on both triples is the tighter constraint.

Conversely, imposing $Q(c,b,t) = 2/3$ exactly on the self-consistent
$m_c^{\text{seed}}$ and $m_b^{v_0}$ predicts
$m_t = 171.3 \pm 1.5\GeV$ (where the uncertainty comes from
$\delta m_s = 0.8\MeV$), within $0.9\%$ of the PDG pole mass
$m_t = 172.76 \pm 0.30\GeV$.  The entire third-generation mass
spectrum then follows from a single input~$m_s$; using
$m_s = 94.2\MeV$ (within $1\sigma$ of PDG) gives $m_t$ exactly.

\paragraph{Summary of predictions.}
Table~\ref{tab:predictions} collects the quantitative predictions.
\begin{table}[ht]
\centering
\small
\begin{tabular}{llll}
\toprule
Quantity & Predicted & PDG 2024 & Source \\
\midrule
$m_c$ & $1{,}301\MeV$ & $1{,}270 \pm 20$ ($1.6\sigma$)
  & Seed: $gv/m = \sqrt{3}$ \\
$m_b$ & $4{,}177\MeV$ & $4{,}180 \pm 30$ ($0.1\sigma$)
  & Bion K\"ahler \\
$m_t$ & $171.3 \pm 1.5\GeV$ & $172.76 \pm 0.30$
  & $Q(c,b,t) = 2/3$ \\
$\sin^2\theta_W$ & $0.22310$ & $0.22320 \pm 0.00026$ ($0.4\sigma$)
  & Casimir equation \\
$V_{us}$ & $0.224$ & $0.2243 \pm 0.0005$ ($1.4\sigma$)
  & Oakes relation \\
$M_0$ (lepton) & $\fpi^2/27 = 314.0\MeV$ & $313.8\MeV$ (fitted)
  & Scale relation \\
$\tan\beta$ & $1$ & --- (not directly measured)
  & Mixed-triple Koide \\
\bottomrule
\end{tabular}
\caption{Predictions from the sBootstrap Lagrangian.
The quark masses use $m_s = 93.4\MeV$ as input;
the bion $m_b$ uses PDG $m_c$.
The Casimir and Oakes predictions are independent
of~$m_s$.}
\label{tab:predictions}
\end{table}

%======================================================================
\section{The Electron Mass and the Lepton Seed}
\label{sec:electron}
%======================================================================

The charged lepton triple $(e, \mu, \tau)$ is the ``almost-seed'': 
$\zch_e = 0.715\MeV^{1/2}$ is small but nonzero. The electron's 
charge is near zero because of cancellation in $\zch_e = z_0 + z_e$:
\begin{equation}
v_0 = 17.716, \quad z_e = -17.001, \quad 
\zch_e = 17.716 - 17.001 = 0.715\MeV^{1/2}\,.
\end{equation}
A 96\% cancellation. The mass $m_e = 0.715^2 = 0.511\MeV$ is the 
square of something already small.

The near-vanishing of $m_e$ is a structural property of the Koide 
embedding, not fine-tuning: the energy balance allows perturbations 
as large as $z_0$. Every Koide triple has one member with 
$|\zch|/v_0$ small:

\begin{center}
\begin{tabular}{lccl}
\toprule
Triple & Lightest & $|\zch|/v_0$ & Character \\
\midrule
$(e, \mu, \tau)$ & $e$ & 0.040 & Almost-seed \\
$(-\pi, D_s, B)$ & $\pi$ (neg.) & 0.337 & Full triple \\
$(c, b, t)$ & $c$ & 0.236 & Full triple \\
$(-s, c, b)$ & $s$ (neg.) & 0.320 & Full triple \\
\bottomrule
\end{tabular}
\end{center}

The electron is by far the most extreme case, consistent with the 
lepton triple being nearest to its seed.

%======================================================================
\section{Symmetry Restoration}
\label{sec:restoration}
%======================================================================

The broken spectrum reassembles as symmetry-breaking effects are 
removed:

\begin{enumerate}
\item \textbf{Turn off QED:} Electromagnetic splittings vanish 
($\pi^\pm = \pi^0$, $K^\pm = K^0$). D-term effects disappear.
\item \textbf{Set $m_u = m_d$:} Exact isospin.
\item \textbf{Set $m_s = m_d$:} $SU(3)_V$ restored, 
$\pi = K = \eta$. Koide splitting $v_1 \to 0$; lepton partners 
degenerate.
\item \textbf{Send $m_q \to 0$:} Chiral limit, all meson charges 
$\to 0$. If sBootstrap SUSY restores, all light lepton masses 
vanish too. This is the seed configuration: $(0, 0, 0)$.
\end{enumerate}

The seed triples are the ``penultimate'' stage: one member has 
reached zero, but the other two retain the $2 - \sqrt{3}$ ratio. 
Full triples are the final broken state. The lepton triple, 
trapped between seed and bloom, reveals the hierarchy: 
chiral breaking $\to$ seed $\to$ bloom $\to$ full spectrum.

The direction of the flow also matters: approaching the chiral 
limit, the pion charge rises from $\zch = -11.81$ toward zero 
from below. It crosses zero exactly in the chiral limit. This 
means the pion was ``born negative''---its charge sign is set by 
the direction of chiral symmetry breaking, and the seed 
$(0, \pi, D_s)$ is the configuration at the moment of crossing.

%======================================================================
%======================================================================
\section{Open Problems}
\label{sec:open}
%======================================================================

\subsection{The diquark sector}
The Pauli theorem (\S\ref{subsec:pauli}) proves that no $J=0$, 
color-$\bar{3}$ diquark can have symmetric flavor, closing all 
field-theoretic escape routes ($L=0$, P-wave, or higher). The 
sBootstrap's $\mathbf{15}$ diquark must therefore be non-composite: 
either an algebraic image of the Miyazawa supercharge
(\S\ref{subsec:baryon_open}) or a fundamental $SO(32)$ string state.
However, the mixed anticommutator
$\{Q_\alpha, \mathcal{Q}_{[ij],\beta}\} = \epsilon_{\alpha\beta}\,
\mathcal{S}_{[ij]}$ produces a diquark generator in
$\overline{\mathbf{10}}$ (antisymmetric), not the symmetric
$\mathbf{15}$; thus the Miyazawa route alone does not resolve the
$\mathbf{15}$~vs.~$\overline{\mathbf{10}}$ conflict.
The exotic $+4/3$-charged
diquarks of the $\mathbf{15}$ would be a testable consequence:
if their mass is set by the SUSY-breaking scale ($\sim f_\pi$), they
would be sub-GeV and accessible at fixed-target experiments; if by
$\Lambda_{\text{QCD}}$ or higher, LHC pair-production searches for
exotic charges apply. No mass prediction can be made without a
concrete realization of the $\mathbf{15}$.

\subsection{What determines the Koide phase?}
The energy balance constrains but does not uniquely fix $\delta$.
The O'Raifeartaigh realization (\S\ref{subsec:fterm}) fixes the
seed ratio $gv/m = \sqrt{3}$ but leaves $\delta$ undetermined---the
phase controls how the seed blooms into a full triple. The bloom
is a phase rotation $\delta: 3\pi/4 \to \delta_{\text{phys}}$ at
fixed $v_0$, but explicit R-breaking perturbations push $Q$ away
from $2/3$ at any nonzero strength, making the seed an unstable
fixed point. The bloom must therefore be nonperturbative.

The lepton phase satisfies $\delta_\ell \bmod 2\pi/3 = 2/9$ to
$33\,\text{ppm}$ ($0.9\sigma$ in $m_\tau$; $3.3\sigma$ after
look-elsewhere over fractions with denominator $\leq 20$). Since $2/9$
is not a rational multiple of $\pi$, this is not a geometric angle.
It is, however, $2/9 = 2/(3\times 3)$, which points to a $\mathbb{Z}_9$
discrete symmetry in the Koide modulus space. The $\mathbb{Z}_3$
mass-permutation symmetry ($\delta \to \delta + 2\pi/3$) is the
unique subgroup; the quotient $\mathbb{Z}_9/\mathbb{Z}_3$ has three
cosets, producing three inequivalent mass spectra from the same $M_0$.
The simplest potential with $\mathbb{Z}_9$ symmetry,
$V \supset -\cos\bigl(9(\delta - 2/9)\bigr)$, has the lepton value
as one of its nine minima, with the $\mathbb{Z}_3$-equivalent minima
at $\delta = 2/9 + 2k\pi/3$.

A dynamical origin for $\mathbb{Z}_9$ exists in SQCD: for $N_f = N_c = 3$,
the diagonal $U(1)_{R+2A}$ (R-symmetry plus twice the axial symmetry)
has anomaly coefficient $6N_c = 18 = 2 \times 9$. Instantons break
$U(1)_{R+2A} \to \mathbb{Z}_{18}$, which contains $\mathbb{Z}_9$
as a subgroup. The meson field $M$ carries $\mathbb{Z}_{18}$ charge 4
(from two quark zero modes at charge 2 each), so $\det M$ carries
charge~12.  The lowest-order holomorphic $\mathbb{Z}_{18}$-singlet
beyond the Seiberg constraint is the three-instanton operator
\begin{equation}
  \delta W_{3\text{-inst}}
    = c_3\,\frac{(\det M)^3}{\Lambda^{18}}\,,
\end{equation}
since $(\det M)^3$ has $\mathbb{Z}_{18}$ charge $36 = 2\times 18$.
Near the SUSY vacuum $\det M = \Lambda^6$, parametrizing
$\det M = \Lambda^6\,e^{3i\delta_M}$ yields
$\delta W \propto e^{9i\delta}$, and the cross-term with the
tree-level superpotential gives
$V_{\text{eff}} \supset -V_0\cos\bigl(9(\delta - \delta_0)\bigr)$
with $\delta_0 = \arg(c_3)/9$.
This is exactly the $\cos(9\delta)$ potential identified above.
The one-instanton vertex ($\det M$, charge 12) is already captured
by the Seiberg constraint; the two-instanton vertex ($(\det M)^2$,
charge 24, not a $\mathbb{Z}_{18}$ singlet) is forbidden.
The three-instanton correction does not affect the meson mass
spectrum: the $F$-term equations define an effective Lagrange
multiplier $X_{\text{eff}} = X + 3c_3\Lambda^{12}/\Lambda^{18}$,
but the constraint $\det M = \Lambda^6$ and the meson VEVs
$M_j = C/m_j$ are unmodified. The correction acts only on the
phase structure (through $\cos(9\delta)$), not on the magnitudes.

At strong coupling ($\Lambda \sim 300\MeV$), the coefficient
$c_3$ is $O(1)$ and the potential strength is
$V_0^{1/4} \sim 65\text{--}85\MeV$, comparable to
$f_\pi = 92\MeV$---this is not a parametrically suppressed effect.
The same $\cos(9\delta)$ structure also arises from magnetic
bion contributions on $\mathbb{R}^3 \times S^1$, where the
$\mathbb{Z}_3$-symmetric combination of the SU(3) affine roots
maps to $\cos(9\delta)$ in the meson phase variable.

The embedding is specific to $N_c = 3$: the condition
$N_c^2 \mid 6N_c$ requires $N_c \mid 6$, which holds only for
$N_c = 1, 2, 3, 6$.

The quark Koide phases, by contrast, do \emph{not} sit at
$\mathbb{Z}_9$ minima: the $(-s,c,b)$ and $(c,b,t)$ phases
are displaced by $39\%$ and $22\%$ of the $\mathbb{Z}_9$
spacing, respectively.  This suggests that the lepton and quark
phases are selected by different dynamical mechanisms---the
instanton potential for leptons, and possibly the bloom mechanism
or ISS dynamics for quarks.

The Coleman--Weinberg potential of the minimal O'Raifeartaigh model
is flat in $\delta$ at one loop---the pseudo-modulus potential depends
only on the modulus $|\Phi_0|$, not on the Koide phase.  The
three-instanton potential above provides the leading
$\delta$-dependent term that lifts this flat direction.
The \emph{modulus} direction $|\Phi_0|$ is a separate question:
the CW potential stabilizes $|\langle\Phi_0\rangle|$ at zero,
but the Koide seed requires $\langle X\rangle/\mu = \sqrt{3}$
(\S\ref{subsec:superpotential}).
A negative quartic K\"ahler correction $c\,|X|^4/\Lambda_K^2$
with $c = -\Lambda_K^2/(12\mu^2)$ pins the pseudo-modulus at
the K\"ahler pole $v_{\text{pole}} = \sqrt{3}\,\mu$ exactly
(an algebraic result, not numerical tuning;
see \S\ref{subsec:superpotential}).
This coefficient is \emph{not} radiatively generated: the
one-loop CW quartic in the O'Raifeartaigh model has $c_4 = -20/3$
at $gf/m^2 = 1$ (exact), which is orders of magnitude too large;
the ISS one-loop potential produces a positive quartic ($c_{\text{eff}} = +0.014$),
the wrong sign entirely. The K\"ahler pole mechanism is therefore a
tree-level statement about the effective K\"ahler potential,
not a consequence of one-loop dynamics.  However, magnetic bions
on $\mathbb{R}^3 \times S^1$ generate a quartic K\"ahler correction
of the correct sign.
The bion-induced coefficient
$c_{\text{bion}} = -3\pi^2\,e^{-4\pi/(N_c^2\alpha_s)}\,M^2/\alpha_s^2$
(where $M^2 = [\prod_f(m_f/\Lambda)]^2$ is the zero-mode mass factor)
equals $-1/12$ at $\alpha_s = 0.143$ ($S_{\text{bion}} = 9.76$)
in the ISS model with heavy magnetic quarks ($M^2 = 1$).
With physical light-quark masses, zero-mode suppression reduces
$|c_{\text{bion}}|$ by five orders of magnitude; the ISS magnetic
dual---where the pseudo-modulus lives---is the appropriate
setting for the bion calculation.
The three-instanton operator $(\det M)^3$ shifts the Seiberg
Lagrange multiplier~$X$ (dimension $[\text{mass}]^{-3}$ in the
confining phase) but this is a different object from the ISS
pseudo-modulus (dimension $[\text{mass}]^{+1}$ in the magnetic
dual).

More broadly, the ISS Coleman--Weinberg potential does \emph{not}
transmit the Koide condition from the UV quark masses to the IR
meson spectrum.  In the $N_f = 4$ ISS model, the one-loop CW
masses for the broken mesons $\chi_a$ are nearly degenerate
($m_{\text{CW}} \approx 27\MeV$ for all four flavors), because
the CW mass depends on $m_a$ through the affine map
$F_a = h\mu^2 - m_a$, and all quark masses satisfy
$m_a \ll h\mu^2$.  The resulting spectrum gives
$Q \approx 1/3$ (the degenerate limit), regardless of the UV
value.  This confirms that flavor-sensitive mass structure must
arise from non-perturbative effects---the bion K\"ahler
correction (\S\ref{subsec:fterm}), which directly involves
$\sqrt{m_k}$, rather than the one-loop effective potential.

\subsection{Connecting the three tuples}
The unified superpotential (\S\ref{subsec:superpotential}) connects
the light and heavy sectors through a layered structure:
the seed Koide fixes $m_c/m_s$ (from $gv/m = \sqrt{3}$),
the $v_0$-doubling fixes $m_b$ (from the K\"ahler sector),
and the $(c,b,t)$ Koide serves as an overdetermination test.
After these constraints, only two to three free parameters remain:
$m_u$, $m_d$ (connected to the Cabibbo angle), and the overall
scale $m_s$. The $v_0$-doubling is non-holomorphic and lives in the
K\"ahler sector: as shown in \S\ref{subsec:fterm}, the bion
K\"ahler correction $K_{\text{bion}} \sim |\!\sum s_k\sqrt{\hat m_k}|^2$
generates an effective potential whose minimum enforces
$\sqrt{m_b} = 3\sqrt{m_s} + \sqrt{m_c}$.
The Hessian has rank~1, constraining exactly the $v_0$-doubling
combination while leaving $m_s$ and $m_c$ to the Koide seed.
The extension from $\mathbb{R}^3 \times S^1$ to
$\mathbb{R}^4$ relies on \"Unsal's continuity conjecture.

\subsection{The baryon sector}
\label{subsec:baryon_open}
In Seiberg's SQCD with $N_f = N_c = 3$, baryons are part of the
low-energy spectrum. Their SUSY partners are ``baryonic fermions''
with no known candidates among observed particles.

The meson--baryon connection requires a Miyazawa supercharge
\cite{Miyazawa, CattoGursey}, which cannot lie in the
fundamental~$\mathbf{5}$: the product
$\mathbf{5}\otimes(\mathbf{1}\oplus\mathbf{24})$ does not
contain~$\overline{\mathbf{10}}$, so a single-index generator
cannot map mesons to baryons.  The minimal choice is
\begin{equation*}
\mathcal{Q}_{[pq],\alpha}
  \;\in\; \overline{\mathbf{10}}\;\text{of}\;SU(5)_f\,,
  \qquad \Delta B = +1\,,
\end{equation*}
with transformation law
$[\mathcal{Q}_{[pq]}, M^i{}_j]
  = \delta^i_{[p}\,\tilde\psi_{B,\,q]j}
  - \tfrac{1}{5}\delta^i_j\,\tilde\psi_{B,\,pq}$,
where $\tilde\psi_{B,\,lm}
  = \tfrac{1}{6}\epsilon_{lmijk}\,\psi_B^{ijk}$ is the dual
baryonino.  The anticommutator decomposes as
\begin{equation*}
\overline{\mathbf{10}}\otimes\mathbf{10}
  = \mathbf{1}\oplus\mathbf{24}\oplus\mathbf{75}\,,
\end{equation*}
giving $\{\mathcal{Q},\bar{\mathcal{Q}}\}
  \supset \sigma^\mu P_\mu + T^A + Z$, where $T^A$ are the
$SU(5)_f$ generators and $Z$ are central charges encoding the
meson--baryon mass splitting $m_B^2 - m_M^2 = \langle Z\rangle
\sim\Lambda_{\text{QCD}}^2$.  This Miyazawa SUSY evades the
Haag--\L{}opusza\'nski--Sohnius theorem because it is an
approximate symmetry broken at scale~$\Lambda_{\text{QCD}}$,
analogous to $SU(6)$ spin--flavor symmetry in the quark model.
The $\mathbf{75}$ channel on the right-hand side
signals that the full bosonic subalgebra extends beyond
$SU(5)_f$; its interpretation remains open.
The $\overline{\mathbf{10}}$ representation of the Miyazawa
supercharge is the same $SU(5)_f$ multiplet that appears in the
$SO(32)$ adjoint decomposition (Table~\ref{tab:496}), suggesting
that the baryon--meson connection and the string embedding may
share a common algebraic origin.

\subsection{The strong CP problem}
The principle ``Nature abhors chirality'' predicts $\bar{\theta} = 0$ 
but provides no mechanism.

\subsection{The kaon--muon puzzle}
The Lagrangian predicts a universal soft splitting
$m^2_{\text{meson}} - m^2_{\text{lepton}} = \fpi^2$ across the
meson matrix. For the pion--muon pair, this works:
$m_\pi^2 - m_\mu^2 \approx \fpi^2$ to $2.2\%$. But the
tree-level formula $m^2_M = |m_i + m_j|^2 + \fpi^2$ badly
underpredicts heavier meson masses: it gives $m_K \approx 133\MeV$
(physical: $498\MeV$) and $m_{D_s} \approx 1{,}367\MeV$
(physical: $1{,}968\MeV$). QCD confinement generates the bulk
of these masses independently of the soft term. The universal
splitting applies to the \emph{soft} contribution, not to the
total mass; it is visible only for $(\pi, \mu)$,
where the soft term dominates over $|m_u + m_d|^2$.

\subsection{The Casimir eigenvalue equation}
The eigenvalue equation $x^2 + C_2 x - C_2 = 0$ produces the
on-shell weak mixing angle to $0.4\sigma$, but three issues
remain open: (i)~the equation is an ansatz, not derived from a
Lagrangian; (ii)~the match uses the on-shell renormalization
scheme, and no argument selects this scheme over $\overline{\text{MS}}$;
(iii)~the negative spin-$\tfrac{1}{2}$ root gives $|M| = 122.4\GeV$,
tantalizingly close to $M_H = 125.25\GeV$ but with a spin mismatch.

\subsection{The Higgs mass at $\tan\beta = 1$}
At $\tan\beta = 1$, the $D$-term quartic
$(g^2+g'^2)(|H_u|^2 - |H_d|^2)^2/8$ vanishes on the
$D$-flat direction $H_u = H_d$.  The tree-level lightest
CP-even Higgs mass is $m_h = m_Z|\cos 2\beta| = 0$---a known
MSSM result. The meson Yukawa F-terms contribute only mass
terms ($y_j^2|H|^2$), not quartics.
In the NMSSM, a singlet coupling
$\lambda\,S\,H_u\!\cdot\!H_d$ generates an extra quartic
$\lambda^2|H_u\!\cdot\!H_d|^2$; at $\sin 2\beta = 1$ (i.e.,
$\tan\beta = 1$), $m_h^2 = \lambda^2 v^2/2$, and
$\lambda = 0.72$ gives $m_h = 125\GeV$ at tree level.
The Seiberg Lagrange multiplier~$X$ is a gauge-singlet chiral
superfield already present in the Lagrangian; adding the coupling
$\lambda\,X\,H_u\!\cdot\!H_d$ to the superpotential is
the minimal extension.

\subsection{The $SO(32)$ projection}
The adjoint $\mathbf{496}$ contains the sBootstrap
representations: $\mathbf{24}$ in the neutral sector and
$\mathbf{15} \oplus \overline{\mathbf{15}}$ in the matter
sector at $|Q_A| = 2$. The remaining 442 states must be
projected out or made massive. No projection mechanism has been
identified. The alternative anomaly-free gauge group
$E_8 \times E_8$ is blocked by the $\mathbf{15}$ vs.\
$\mathbf{10}$ obstruction (\S\ref{subsec:so32_open}),
constraining the string embedding to the Type~I / $SO(32)$
corner of the landscape.

\subsection{Generation uniqueness}
The Diophantine system uniquely selects $N = 3$ generations,
but the third equation ($r^2 + s^2 - 1 = 4N$) has the weakest
physical motivation. A dynamical derivation of the counting
constraints would strengthen the argument considerably.

%======================================================================
\section{Epistemological classification}
\label{sec:epistemology}
%======================================================================

The following table classifies the paper's claims by their logical
status.  ``Derived'' means following from stated assumptions by
algebra or field theory; ``observed'' means an empirical pattern;
``assumed'' means an input without derivation.

\begin{center}
\small
\begin{tabular}{lll}
\toprule
Claim & Status & Comment \\
\midrule
$Q = 2/3$ iff $\langle z_k^2\rangle = z_0^2$ & Derived & Algebra (\S\ref{sec:koide_algebra}) \\
Seed ratio $= 2-\sqrt{3}$ & Derived & From $Q = 2/3$ + one zero \\
O'R seed at $gv/m = \sqrt{3}$ & Derived & Fermion eigenvalues \\
K\"ahler stabilization: $c = -1/12$ & Derived & Pole condition (\S\ref{subsec:superpotential}) \\
$\langle X\rangle/\mu = \sqrt{3}$ in magnetic dual & Derived & ISS dictionary \\
Pauli theorem (no symmetric $J{=}0$ diquark) & Derived & Angular momentum \\
Diophantine: $(N,r,s) = (3,3,2)$ unique & Derived & Number theory \\
Cyclic echo ($t \to c \to s \to \text{stall} \to c$) & Derived & Algebraic necessity \\
Seed self-dual: $Q(a,b) = Q(1/a,1/b)$ & Derived & Identity for two masses \\
Cartan angle $\theta = \pi/6 - \delta$ & Derived & Exact trigonometric identity \\
F-terms flavour-universal: $F = C/v_{\text{EW}}$ & Derived & Seesaw cancellation \\
\midrule
Lepton Koide $Q = 2/3$ to $0.001\%$ & Observed & $2.8\sigma$ after LEE \\
Meson triple $(-\pi,D_s,B)$ to $0.10\%$ & Observed & $1.1\sigma$ after LEE \\
$v_0$-doubling: $v_0^{(f)}/v_0^{(s)} = 2.0005$ ($0.025\%$); $m_b$ to $0.07\%$ & Observed & Derived via bion K\"ahler \\
$\delta_\ell \bmod 2\pi/3 = 2/9$ to 33~ppm & Observed & $3.3\sigma$ after LEE \\
Casimir $R = 0.22310$ vs $\sin^2\theta_W$ & Observed & Scheme-dependent \\
Dual Koide $Q(1/m_d,1/m_s,1/m_b) = 0.665$ & Observed & $0.2\%$ from $2/3$ \\
$M_0 \approx \fpi^2/27$ & Observed & $0.04\%$; connects lepton scale to $\fpi$ \\
$\tan\beta = 1$ & Derived & Mixed triples + $SU(5)_f$ \\
\midrule
$m_\pi^2 - m_\mu^2 = f_\pi^2$ (VA relation) & Assumed & 2.2\% match \\
$f_{\text{VA}} = f_\pi$ identification & Assumed & Central ansatz \\
SU(5) flavor symmetry & Assumed & 5 light quarks \\
$m_b = 4{,}177 \pm 40\MeV$ from bion minimum & Derived & $0.06\sigma$; requires \"Unsal conjecture \\
$N_f = 4$ unique in ISS window & Derived & Only integer in $(3, 4.5)$ \\
ISS vacuum cosmologically stable & Derived & $S_{\text{bounce}} \gg 1$ for $h \sim O(1)$ \\
ISS CW does not transmit Koide & Derived & $Q_{\text{CW}} \to 1/3$ (affine map) \\
$c = -1/12$ from bion K\"ahler & Derived & At $\alpha_s = 0.143$ (ISS) \\
$m_h = 0$ at tree-level ($\tan\beta = 1$) & Derived & NMSSM singlet: $\lambda = 0.72$ \\
$\cos(9\delta)$ from $(\det M)^3/\Lambda^{18}$ & Derived & $\mathbb{Z}_{18}$ charge counting \\
$\mathcal{Q}_{[ij]} \in \overline{\mathbf{10}}$ of $SU(5)_f$ & Derived & Miyazawa algebra \\
Casimir structural constraints & Assumed & Reverse-engineered \\
\bottomrule
\end{tabular}
\end{center}

%======================================================================
\section{Conclusions}
\label{sec:conclusions}
%======================================================================

Seven principal results:

\paragraph{1.\ The Koide formula is energy equipartition.}
$Q = 2/3$ is $\langle z_k^2\rangle = z_0^2$: perturbation energies 
match the base scale. The signed-charge framework resolves the 
electron mass puzzle (near-cancellation, not fine-tuning) and reveals 
that mesons satisfy Koide with $(-\pi, D_s, B)$ to $0.1\%$.

\paragraph{2.\ Seed triples set the stage for SUSY breaking.}
The quark seed $(0,s,c)$ and meson seed $(0,\pi,D_s)$ share the
universal ratio $2 - \sqrt{3}$ and operate at $v_0^2 \sim
230$--$350\MeV$, close to the lepton triple's $v_0^2 = 314\MeV
\approx \fpi^2/27$
(recall that $v_0^2$ carries dimensions of MeV, not MeV$^2$,
because $v_0 = z_0$ is a charge with dimensions MeV$^{1/2}$). 
The massless member is the necessary condition for F-term 
(O'Raifeartaigh) SUSY breaking. Upon breaking, seeds bloom into 
full triples $(-s,c,b)$ and $(-\pi,D_s,B)$. The lepton triple 
is the ``almost-seed'' that has barely bloomed.

\paragraph{3.\ The seed condition is one superpotential constraint,
stabilized by a K\"ahler pole.}
The O'Raifeartaigh model with superpotential
$W = f\Phi_0 + m\Phi_1\Phi_2 + g\Phi_0\Phi_1^2$ produces a
fermionic spectrum $(0,\, m_-,\, m_+)$ that is an exact Koide
seed triple when $gv/m = \sqrt{3}$. The one-loop CW potential
alone stabilizes the pseudo-modulus at $v = 0$; a negative
quartic K\"ahler correction $c\,|\Phi_0|^4/\Lambda_K^2$ with
$c = -\Lambda_K^2/(12m^2)$ pins it at the K\"ahler pole
$v = \sqrt{3}\,m/g$ exactly. The Koide seed emerges at
$M_0 = 2m/3$.  The coefficient $c = -1/12$ is not radiatively
generated but arises from magnetic bions on
$\mathbb{R}^3 \times S^1$ at $\alpha_s = 0.143$
in the ISS magnetic dual.

\paragraph{4.\ Two breaking mechanisms plus a heavy-sector origin.}
F-term breaking (via the massless seed member) generates the light 
sector; the Volkov--Akulov relation $f_{\text{VA}} = f_\pi$ is its 
low-energy realization, not a separate mechanism. D-term breaking 
($U(1)_{\text{em}}$ FI parameter) provides electromagnetic 
splittings as a perturbative correction. The heavy triple 
$(c,b,t)$ stands apart: no seed, no meson shadow, scale far above 
$\Lambda_{\text{QCD}}$, possibly connected to electroweak breaking.

\paragraph{5.\ The Diophantine counting uniquely selects $N=3$.}
Three equations on the scalar content---$rs = 2N$, $r(r+1)/2 = 2N$,
$r^2+s^2-1 = 4N$---have the unique positive-integer solution
$(N,r,s) = (3,3,2)$, which maps to the $SU(5)$ representations
$\mathbf{24} \oplus \mathbf{15} \oplus \overline{\mathbf{15}}$.
This combination is anomaly-free and asymptotically free.

\paragraph{6.\ An algebraic electroweak ratio.}
The eigenvalue equation $x^2 + s(s+1)\,x - s(s+1) = 0$, with
$s$ the spin Casimir, predicts
$R = 1 - x_+(1/2)/x_+(1) = 0.22310$, matching the on-shell
$\sin^2\theta_W = 0.22320$ to $0.4\sigma$ and giving
$M_W = 80.374\GeV$ ($0.39\sigma$ from PDG). The equation is
pinned by three structural constraints; $R$ follows as a
consequence, not a free parameter. The dynamical origin of the
equation and the choice of renormalization scheme remain open.

\paragraph{7.\ An explicit Lagrangian with two to three free parameters.}
The full $\mathcal{N}=1$ Lagrangian
(\S\ref{subsec:superpotential}) comprises a K\"ahler potential
with a negative quartic correction $c = -1/12$ (stabilizing
the pseudo-modulus at $gv/m = \sqrt{3}$) and a bion correction
$K_{\text{bion}} \sim |\!\sum s_k\sqrt{\hat m_k}|^2$,
a superpotential combining the Seiberg constraint
with Yukawa couplings (using separate $H_u$, $H_d$ doublets),
$D$-terms with a Fayet--Iliopoulos
parameter, and soft breaking from a spurion
$\langle F\rangle = \fpi^2$.
The mixed up/down-type Koide triples require $\tan\beta = 1$
for the mass-level condition to equal a Yukawa-level
condition---a prediction independently motivated by $SU(5)_f$.
The seed Koide fixes the
pseudo-modulus VEV $\langle X\rangle = \sqrt{3}\,\mu$ (predicting $m_c/m_s$); the
bion K\"ahler correction fixes $m_b$ through the
$v_0$-doubling condition, with Hessian rank~1 constraining
exactly one mass combination.  The self-consistent chain from $m_s$ alone predicts
$m_c = 1{,}301\MeV$, $m_b = 4{,}177\MeV$, and---if
$Q(c,b,t) = 2/3$ is imposed---$m_t = 171.3 \pm 1.5\GeV$
(Table~\ref{tab:predictions}).
The lepton Koide scale satisfies $M_0 \approx \fpi^2/27$
to $0.3\%$, tying lepton masses to the pion decay constant.
After these
constraints, only $m_u$, $m_d$, and $m_s$ remain free; $m_u$ and
$m_d$ are connected to the Cabibbo angle through the Oakes
relation.  The $SO(32)$ adjoint provides the UV completion,
with the $E_8 \times E_8$ alternative blocked by the
$\mathbf{15}$-vs.-$\mathbf{10}$ obstruction.

The diquark sector, required to complete the SUSY multiplet count,
fills the symmetric $\mathbf{15} + \overline{\mathbf{15}}$ of
$SU(5)_f$---the representation produced naturally by the SUSY
generator acting on quarks and by the $SO(32)$ decomposition. A
kinematic theorem (\S\ref{subsec:pauli}) shows that \emph{no}
$J=0$ color-$\bar{3}$ diquark can have symmetric flavor, regardless
of orbital angular momentum: the sBootstrap diquark in the
$\mathbf{15}$ cannot be a $qq$ composite. The Miyazawa supercharge
$\mathcal{Q}_{[ij]} \in \overline{\mathbf{10}}$
(\S\ref{subsec:baryon_open}) connects mesons to baryons but
generates diquark operators in $\overline{\mathbf{10}}$
(antisymmetric), not $\mathbf{15}$ (symmetric).  The remaining
candidate is a fundamental string-theoretic state from the
$SO(32)$ adjoint.

The organizing principle is that each sector of the spectrum begins
with a massless member (a seed), which is the necessary condition
for spontaneous SUSY breaking. The breaking lifts the zero,
generates the full triple, and the energy-balance condition
$\langle z_k^2\rangle = z_0^2$ constrains the initial conditions.
The bloom from seed to full triple is a phase rotation
$\delta: 3\pi/4 \to \delta_{\text{phys}}$ in the Koide
parametrization, at fixed Koide scale $v_0$; it is not a
continuous deformation of the O'Raifeartaigh spectrum (which would
break $Q = 2/3$) but a nonperturbative rearrangement.  The lepton
phase is selected by a three-instanton potential
$V \propto -\cos(9\delta)$, the lowest $\mathbb{Z}_{18}$-invariant
harmonic in $N_f = N_c = 3$ SQCD; the quark phases are not at
$\mathbb{Z}_9$ minima, suggesting a different selection mechanism.
The unified superpotential has a layered structure:
(i)~the Seiberg effective superpotential governs the light sector,
with SUSY breaking forced by the quantum constraint when $m_u = m_d = 0$;
(ii)~the O'Raifeartaigh condition $gv/m = \sqrt{3}$ (equivalently
$\langle X\rangle = \sqrt{3}\,\mu$ in the magnetic dual) fixes the seed Koide ratio,
predicting $m_c$ from $m_s$;
(iii)~the bion K\"ahler correction
$|\!\sum s_k\sqrt{\hat m_k}|^2$ predicts
$m_b = 4{,}177 \pm 40\MeV$ via $v_0$-doubling;
(iv)~the Yukawa sector governs the $(c,b,t)$ triple as an
overdetermination test.
After these constraints, only two to three free parameters remain:
$m_u$, $m_d$ (connected to the Cabibbo angle through the Oakes
relation), and the overall scale $m_s$. All quark Koide triples
obligatorily mix up-type and down-type mass eigenstates.
The energy-balance condition is a boundary condition on the
superpotential, not an RG attractor; the Koide ratios are naturally
protected against $D$-term perturbations because $D$-terms shift
scalar masses but leave fermion masses untouched.

\begin{thebibliography}{99}
\bibitem{Rivero2005}
A.~Rivero, ``An interpretation of scalars in SO(32),''
Eur.\ Phys.\ J.\ C \textbf{84}, 1058 (2024),
arXiv:2407.05397 [hep-ph]; earlier versions: viXra:0503.001 (2005),
viXra:0809.001 (2008).

\bibitem{Rivero2012}
A.~Rivero, ``A new Koide tuple: strange, charm, bottom,'' 
arXiv:1111.7232 [hep-ph] (2011).

\bibitem{Koide1983}
Y.~Koide, ``New viewpoint of quark and lepton mass hierarchy,'' 
Phys.\ Rev.\ D \textbf{28}, 252 (1983).

\bibitem{CattoGursey}
S.~Catto and F.~G\"ursey, ``Algebraic treatment of effective 
supersymmetry,'' Nuovo Cim.\ A \textbf{86}, 201 (1985).

\bibitem{Miyazawa}
H.~Miyazawa, ``Baryon Number Changing Currents,'' 
Prog.\ Theor.\ Phys.\ \textbf{36}, 1266 (1966).

\bibitem{Harari}
H.~Harari, H.~Haut, J.~Weyers, ``Quark masses and Cabibbo 
angles,'' Phys.\ Lett.\ B \textbf{78}, 459 (1978).

\bibitem{Seiberg1995}
N.~Seiberg, ``Electric--magnetic duality in supersymmetric 
non-abelian gauge theories,'' Nucl.\ Phys.\ B \textbf{435}, 129 
(1995).

\bibitem{ORaifeartaigh}
L.~O'Raifeartaigh, ``Spontaneous symmetry breaking for chiral 
scalar superfields,'' Nucl.\ Phys.\ B \textbf{96}, 331 (1975).

\bibitem{FayetIliopoulos}
P.~Fayet and J.~Iliopoulos, ``Spontaneously broken supergauge 
symmetries and Goldstone spinors,'' Phys.\ Lett.\ B \textbf{51}, 
461 (1974).

\bibitem{Witten1983}
E.~Witten, ``Some inequalities among hadron masses,'' 
Phys.\ Rev.\ Lett.\ \textbf{51}, 2351 (1983).

\bibitem{VolkovAkulov1973}
D.~V.~Volkov and V.~P.~Akulov, ``Is the neutrino a Goldstone 
particle?,'' Phys.\ Lett.\ B \textbf{46}, 109 (1973).

\bibitem{GOR1968}
M.~Gell-Mann, R.~J.~Oakes, B.~Renner, ``Behavior of current 
divergences under $SU(3) \times SU(3)$,'' 
Phys.\ Rev.\ \textbf{175}, 2195 (1968).

\bibitem{ISS2006}
K.~Intriligator, N.~Seiberg, D.~Shih, ``Dynamical SUSY breaking 
in meta-stable vacua,'' JHEP \textbf{0604}, 021 (2006),
arXiv:hep-ph/0602239.

\bibitem{DGMLY1967}
T.~Das, G.~S.~Guralnik, V.~S.~Mathur, F.~E.~Low, J.~E.~Young,
``Electromagnetic mass difference of pions,''
Phys.\ Rev.\ Lett.\ \textbf{18}, 759 (1967).

\bibitem{FramptonGlashow1987}
P.~H.~Frampton and S.~L.~Glashow, ``Chiral color: an alternative
to the standard model,'' Phys.\ Lett.\ B \textbf{190}, 157 (1987).

\bibitem{Unsal2008}
M.~\"Unsal, ``Abelian duality, confinement, and chiral symmetry
breaking in QCD(adj),'' Phys.\ Rev.\ Lett.\ \textbf{100}, 032005
(2008), arXiv:0708.1772.

\bibitem{Unsal2012}
M.~\"Unsal, ``Theta dependence, sign problems and topological
interference,'' Phys.\ Rev.\ D \textbf{86}, 105012 (2012),
arXiv:1201.6426.

\bibitem{Rivero2005b}
A.~Rivero, ``The de Vries formula for $\sin^2\theta_W$,''
arXiv:hep-ph/0510068 (2005).

\end{thebibliography}

\end{document}
